\begin{enumerate}
\def\labelenumi{\arabic{enumi}.}
\setcounter{enumi}{10}
\item
  Let \(\mathfrak{g}\) be a Lie algebra and \(\mathfrak{h}\) an ideal of \(\mathfrak{g}\) such that \(\mathcal{D}\mathfrak{h} = \mathfrak{h}\) . Show that \(\mathfrak{h}\) is a characteristic ideal of \(\mathfrak{g}\) .
\item
  Let g be a Lie algebra.
\end{enumerate}

\begin{enumerate}
\def\labelenumi{(\alph{enumi})}
\item
  For a derivation D of g to commute with all the inner derivations of g, it is necessary and sufficient that D map g into its centre.
\item
  Suppose henceforth that \(\mathfrak{g}\) has centre zero, so that \(\mathfrak{g}\) is identified with an ideal of the Lie algebra \(\mathfrak{D}\) of all derivations of \(\mathfrak{g}\) . Show that the centralizer of \(\mathfrak{g}\) in \(\mathfrak{D}\) is zero (use (a)). In particular, the centre of \(\mathfrak{D}\) is zero.
\item
  Show that a derivation \(\Delta\) of \(\mathfrak{D}\) such that \(\Delta(\mathfrak{g}) = \{0\}\) is zero. \(([\Delta(\mathfrak{D}),\mathfrak{g}]\subset \Delta([\mathfrak{D},\mathfrak{g}]) + [\mathfrak{D},\Delta(\mathfrak{g})]\subset \Delta(\mathfrak{g}) = \{0\}\) and hence \(\Delta(\mathfrak{D}) = \{0\}\) by (b).)
\item
  Show that, if further \(\mathfrak{g} = \mathcal{D}\mathfrak{g}\) , every derivation \(\Delta\) of \(\mathfrak{D}\) is inner (use (c) and Exercise 11).
\end{enumerate}

\begin{enumerate}
\def\labelenumi{\arabic{enumi}.}
\setcounter{enumi}{12}
\tightlist
\item
  Let \(\mathfrak{g}\) be a Lie algebra and \(\mathfrak{a}\) and \(\mathfrak{b}\) two submodules of \(\mathfrak{g}\) . For every integer \(i \geqslant 0\) we define the submodules \(\mathfrak{m}_i = \mathfrak{m}_i(\mathfrak{a},\mathfrak{b})\) and \(\mathfrak{n}_i = \mathfrak{n}_i(\mathfrak{a},\mathfrak{b})\) as follows: \(\mathfrak{m}_1 = \mathfrak{b}, \mathfrak{m}_{i+1} = [\mathfrak{m}_i,\mathfrak{b}], \mathfrak{n}_1 = \mathfrak{a}, \mathfrak{n}_{i+1} = [\mathfrak{n}_i,\mathfrak{b}]\) . Show that
\end{enumerate}

\[
[ a, m _ {i} ] \subset n _ {i + 1}
\]

for \(i = 1,2,\ldots\) (argue by induction on \(i\) , observing that

\[
\mathfrak {m} _ {i} ([ a, b ], b) = \mathfrak {m} _ {i} (a, b)
\]

and \(n_{i}([a,b],b)=n_{i+1}(a,b)\) .

¶ 14. Let \(\mathfrak{a}\) be a Lie algebra and \(\mathfrak{b}\) a subalgebra of \(\mathfrak{a}\) . A composition series joining \(\mathfrak{a}\) to \(\mathfrak{b}\) is a decreasing sequence \((\mathfrak{a}_i)_{0\leqslant i\leqslant n}\) of subalgebras of \(\mathfrak{a}\) such that \(\mathfrak{a}_0 = \mathfrak{a}\) , \(\mathfrak{a}_n = \mathfrak{b}\) , \(\mathfrak{a}_{i+1}\) is an ideal of \(\mathfrak{a}_i\) for \(0\leqslant i < n\) . \(\mathfrak{b}\) is called a subinvariant subalgebra of \(\mathfrak{a}\) if there exists a composition series joining \(\mathfrak{a}\) to \(\mathfrak{b}\) .

\begin{enumerate}
\def\labelenumi{(\alph{enumi})}
\item
  Let \(\mathfrak{b}\) be a subinvariant subalgebra of \(\mathfrak{a}\) and \((\mathfrak{a}_i)_{0\leqslant i\leqslant n}\) a composition series joining \(\mathfrak{a}\) to \(\mathfrak{b}\) . Deduce from Exercise 13 that \([\mathcal{C}^{n+p}\mathfrak{b},\mathfrak{a}]\subset\mathcal{C}^{p}\mathfrak{b}\) by observing that \([\mathfrak{a}_i,\mathfrak{b}]\subset\mathfrak{a}_{i+1}\) for \(1\leqslant i<n\) .
\item
  Deduce from (a) that the intersection \(\mathcal{C}^{\infty}\mathfrak{b}\) of the \(\mathcal{C}^{p}\mathfrak{b}\) ( \(p = 0,1,2,\ldots\) ) is an ideal of \(\mathfrak{a}\) .
\item
  Deduce from (b) that a subinvariant subalgebra \(\mathfrak{c}\) of \(\mathfrak{a}\) such that \(\mathfrak{c} = \mathcal{D}\mathfrak{c}\) is an ideal of \(\mathfrak{a}\) .
\item
  When K is a field and \(\dim_{K}a<+\infty\) , show that the intersection \(D^{\infty}b\) of the \(D^{p}b\left(p=0,1,2,\ldots\right)\) is an ideal of a. (Show that this intersection is a subinvariant subalgebra of a and apply (c).)
\end{enumerate}

¶ 15. (a) Let a be a Lie algebra, b a subalgebra of a, c an ideal of b and z the centralizer of c in a. Then {[}b, b{]} ⊂ z. Deduce that b + z is a subalgebra of a in which z is an ideal.

\begin{enumerate}
\def\labelenumi{(\alph{enumi})}
\setcounter{enumi}{1}
\item
  Let g be a Lie algebra, a a subinvariant subalgebra of g and b an ideal of a. If the centralizer of b in a and the centralizer of a in g are zero, the centralizer \(\mathfrak{z}\) of \(\mathfrak{b}\) in \(\mathfrak{g}\) is zero. (Let \(\mathfrak{b} = \mathfrak{a} + \mathfrak{z}\) , which is a subalgebra by \((a)\) ; \(\mathfrak{a}\) is subinvariant in \(\mathfrak{b}\) ; if \(\mathfrak{z} \neq \{0\}\) , then \(\mathfrak{a} \neq \mathfrak{h}\) , hence \(\mathfrak{a}\) is an ideal in \(\mathfrak{a}'\) , with \(\mathfrak{a} \neq \mathfrak{a}' \subset \mathfrak{b}\) ; let \(a' \in \mathfrak{a}'\) , \(a' \notin \mathfrak{a}\) and \(a' = a + z\) , \(a \in \mathfrak{a}\) , \(z \in \mathfrak{z}\) , \(z \neq 0\) ; show that \([z, \mathfrak{a}]\) is contained in \(\mathfrak{a}\) and commutes with \(\mathfrak{b}\) , whence \([z, \mathfrak{a}] = \{0\}\) giving a contradiction.)
\item
  Let \(\mathfrak{g}\) be a Lie algebra with centre zero, \(\mathfrak{D}_1\) the Lie algebra of derivations of \(\mathfrak{g}\) , \(\mathfrak{D}_2\) the Lie algebra of derivations of \(\mathfrak{D}_1, \ldots, \mathfrak{g}\) is identified with an ideal of \(\mathfrak{D}_1\) , \(\mathfrak{D}_1\) with an ideal of \(\mathfrak{D}_2, \ldots\) (cf.~Exercise 12 (b)). Using (b) and Exercise 12 (b), show that the centralizer of \(\mathfrak{g}\) in \(\mathfrak{D}_i\) is zero for all \(i\) . (For a sequel to this exercise, cf.~§ 4, Exercise 15.)
\end{enumerate}

\begin{enumerate}
\def\labelenumi{\arabic{enumi}.}
\setcounter{enumi}{15}
\tightlist
\item
  Let g be a Lie algebra and D the Lie algebra of derivations of g. The identity mapping of D defines a semi-direct product h of D by g called the holomorph of g. g is identified with an ideal and D with a subalgebra of h.
\end{enumerate}

\begin{enumerate}
\def\labelenumi{(\alph{enumi})}
\item
  Show that the centralizer \(\mathfrak{g}^*\) of \(\mathfrak{g}\) in \(\mathfrak{h}\) is the set of \(\operatorname{ad}_{\mathfrak{g}}x - x(x\in \mathfrak{g})\) .
\item
  For every element \(x + d\) of \(\mathfrak{h}\) ( \(x \in \mathfrak{g}, d \in \mathfrak{D}\) ), we write
\end{enumerate}

\[
\theta (x + d) = \mathrm{ad} _ {\mathfrak {g}} x - x + d.
\]

Show that \(\theta\) is an automorphism of \(\mathfrak{h}\) of order 2 which maps \(\mathfrak{g}\) onto \(\mathfrak{g}^*\) , so that \(\mathfrak{h}\) can be considered as the holomorph of \(\mathfrak{g}^*\) .

\begin{enumerate}
\def\labelenumi{(\alph{enumi})}
\setcounter{enumi}{2}
\item
  For every ideal \(\mathfrak{k} \neq \{0\}\) of \(\mathfrak{h}\) , \(\mathfrak{k} \cap (\mathfrak{g} + \mathfrak{g}^*) \neq \{0\}\) . (If \(\mathfrak{k} \cap \mathfrak{g} = \{0\}\) , \(\mathfrak{k} \subset \mathfrak{g}^*\) .)
\item
  If g is commutative, every ideal \(t \neq \{0\}\) of h contains g. If K is a field and \(\dim_{K} g < +\infty\) , deduce that, if \(t \neq g\) , h cannot be the holomorph of t (use (b)).
\end{enumerate}

\begin{enumerate}
\def\labelenumi{\arabic{enumi}.}
\setcounter{enumi}{16}
\tightlist
\item
  Suppose that K is a field. Let T be an indeterminate and \(L = K((T))\) ; L has canonically a valuation (order of the formal power series) which makes it a complete valued field.
\end{enumerate}

\begin{enumerate}
\def\labelenumi{(\alph{enumi})}
\item
  Let \(\mathrm{V}\) be a finite-dimensional vector space over \(\mathbf{K}\) . Then \(\mathrm{V}_{(\mathbb{L})}\) has canonically a complete metrizable topological vector space structure (Topological Vector Spaces, Chapter I, § 2, no. 3, Theorem 2). Let \((x_{p})_{p\in \mathbb{Z}}\) be a family of elements of \(\mathrm{V}\) such that \(x_{p} = 0\) for \(p\) less than some rational integer. Then the series \(\sum_{-\infty}^{+\infty}x_p\mathrm{T}^p\) converges in \(\mathrm{V}_{(\mathrm{L})}\) and every element of \(\mathrm{V}_{(\mathrm{L})}\) can be expressed uniquely in this way. (Take a basis of V.)
\item
  The vector space \(\mathcal{L}_{\mathrm{L}}(\mathrm{V}_{(\mathrm{L})}) = \mathcal{L}_{\mathrm{K}}(\mathrm{V})_{(\mathrm{L})}\) has a canonical complete metrizable topological vector space structure. Every endomorphism of \(\mathrm{V}_{(\mathrm{L})}\) can be expressed uniquely in the form \(\sum_{-\infty}^{+\infty}u_p\mathrm{T}^p\) , where \(u_{p}\in \mathcal{L}_{\mathrm{K}}(\mathrm{V})\) and the \(u_{p}\) are zero for \(p\) less than some rational integer.
\item
  If \(\mathbf{K}\) is of characteristic 0 and \(u\in \mathcal{L}_{\mathrm{K}}(\mathrm{V})\) , we write
\end{enumerate}

\[
e ^ {\mathrm{Tu}} = 1 + \frac {1}{1 !} u \mathrm{T} + \frac {1}{2 !} u ^ {2} \mathrm{T} ^ {2} + \dots \in \mathcal {L} _ {\mathrm{L}} (\mathrm{V} _ {(\mathrm{L})}).
\]

If \(x \in V\) is such that \(ux = 0\) , then \(e^{Tu}.x = x\) .

\begin{enumerate}
\def\labelenumi{(\alph{enumi})}
\setcounter{enumi}{3}
\item
  If W is another finite-dimensional vector space over K and \(v \in \mathcal{L}_{\mathrm{K}}(\mathrm{W})\) , then \(e^{\mathrm{T}(u \otimes 1 + 1 \otimes v)} = e^{\mathrm{T}u} \otimes e^{\mathrm{T}v}\) .
\item
  Deduce from (c) and (d) that, if V has a (not necessarily associative) algebra structure and u is a derivation of V, then \(e^{Tu}\) is an automorphism of the algebra \(V_{(L)}\) .
\end{enumerate}

¶ 18. Let H be a complex Hilbert space, Y and Z continuous Hermitian operators on H and B = {[}Z, Y{]}. Suppose that Y and B are permutable.

\begin{enumerate}
\def\labelenumi{(\alph{enumi})}
\item
  Show that \([Z, Y^n] = nBY^{n-1}\) (argue by induction on \(n\) ).
\item
  Let \(f\) be a continuously differentiable function of a real variable. Deduce from (a) that \([Z, f(Y)] = Bf'(Y)\) .
\item
  Deduce from (b) that \(B = 0\) . (Show that, if \(B \neq 0\) , \(f\) can be chosen such that \(\| f(Y) \| \leqslant 1\) and \(\| Bf'(Y) \|\) is arbitrarily large.)
\item
  Let g be the Lie algebra of continuous operators T on H such that \(T^{*} = -T\) . Deduce from (c) that the conditions \(Y \in g\) , \(Z \in g\) , \([[Z, Y], Y] = 0\) imply \([Z, Y] = 0\) .
\end{enumerate}

\begin{enumerate}
\def\labelenumi{\arabic{enumi}.}
\setcounter{enumi}{18}
\tightlist
\item
  \begin{enumerate}
  \def\labelenumii{(\alph{enumii})}
  \tightlist
  \item
    Let \(\mathbf{L}\) be an associative algebra over \(\mathbf{K}\) . Given \(k\) elements \(x_{1},\ldots ,x_{k}\) of \(\mathbf{L}\) , we write:
  \end{enumerate}
\end{enumerate}

\[
f _ {k} (x _ {1}, \dots , x _ {k}) = \sum_ {\sigma \in \mathfrak {S}_{k}} x _ {\sigma (1)} \dots x _ {\sigma (k)}.
\]

In the algebra \(L \otimes_{K} K[T_{1}, \ldots, T_{k}]\) , where the \(T_{i}\) are indeterminates, \(f_{k}(x_{1}, \ldots, x_{k})\) is the coefficient of \(T_{1} \ldots T_{k}\) in \((x_{1}T_{1} + \cdots + x_{k}T_{k})^{k}\) ; it is also the coefficient of \(T_{1} \ldots T_{k}\) in:

\[
\sum_ {j = 0} ^ {k - 1} \left(x _ {1} \mathrm{T} _ {1} + \dots + x _ {k - 1} \mathrm{T} _ {k - 1}\right) ^ {j} x _ {k} \mathrm{T} _ {k} \left(x _ {1} \mathrm{T} _ {1} + \dots + x _ {k - 1} \mathrm{T} _ {k - 1}\right) ^ {k - j - 1}.
\]

Given two elements \(x, y\) in \(\mathbf{L}\) , we define \(g_{i,k-i}(x,y)\) as the coefficient of \(\mathrm{T}_1^i\mathrm{T}_2^{k - i}\) in \((x\mathrm{T}_1 + y\mathrm{T}_2)^k\) . Show that

\[
i! (k - i)! g _ {i, k - i} (x, y) = f _ {k} \left(t _ {1}, \dots , t _ {k}\right)
\]

where \(t_h = x\) for \(h \leqslant i\) and \(t_h = y\) for \(h \geqslant i + 1\) .

\begin{enumerate}
\def\labelenumi{(\alph{enumi})}
\setcounter{enumi}{1}
\tightlist
\item
  If L is considered as a Lie algebra over K, then in \(\mathcal{L}_{\mathrm{K}}(\mathrm{L})\) :
\end{enumerate}

\[
(\operatorname{ad} x) ^ {n} = \sum_ {i = 0} ^ {n} (- 1) ^ {n - i} \binom {n} {i} \mathrm{L} _ {x} ^ {i} R _ {x} ^ {n - i},
\]

where \(\mathbf{L}_x\) (resp. \(R_{x}\) ) is the multiplication \(y\mapsto xy\) (resp. \(y\mapsto yx\) ).

Deduce that, if K is of characteristic p (p prime), then:

\begin{enumerate}
\def\labelenumi{(\arabic{enumi})}
\tightlist
\item
\end{enumerate}

\[
(\operatorname{ad} x) ^ {p}. y = (\operatorname{ad} x ^ {p}). y\tag{2}
\]

\[
(\operatorname{ad} x) ^ {p - 1}. y = \sum_ {i = 0} ^ {p - 1} x ^ {i} y x ^ {p - 1 - i}.
\]

Deduce from (2) that:

\[
f _ {p - 1} (\operatorname{ad} x _ {1}, \dots , \operatorname{ad} x _ {p - 1}). y = f _ {p} (x _ {1}, \dots , x _ {p - 1}, y)
\]

and:

\[
g _ {i, p - 1 - i} (\mathrm{ad} x, \mathrm{ad} y). y = (p - i) g _ {i, p - i} (x, y).
\]

Conclude that for any two elements x, y of L:

\[
(x + y) ^ {p} = x ^ {p} + y ^ {p} + \Lambda_ {p} (x, y)\tag{3}
\]

where:

\[
\Lambda_ {p} (x, y) = \sum_ {i = 1} ^ {p - 1} (p - i) ^ {- 1} g _ {i, p - 1 - i} (\mathrm{ad} x, \mathrm{ad} y). y
\]

belongs to the Lie subalgebra of L generated by x and y (``Jacobson's formulae'').

\begin{enumerate}
\def\labelenumi{\arabic{enumi}.}
\setcounter{enumi}{19}
\tightlist
\item
  Let g be a Lie algebra over a ring K such that \(pK = \{0\}\) (p prime). A mapping \(x \mapsto x^{[p]}\) of g into itself is called a p-mapping if it satisfies the relations
\end{enumerate}

\[
\begin{array}{r l} & {\operatorname{ad} x ^ {[ p ]} = (\operatorname{ad} x) ^ {p}} \\ & {(\lambda x) ^ {[ p ]} = \lambda^ {p} x ^ {[ p ]}} \\ & {(x + y) ^ {[ p ]} = x ^ {[ p ]} + y ^ {[ p ]} + \Lambda_ {p} (x, y) \quad (\mathrm{cf.Exercise19(b)})} \end{array}
\]

  for \(x, y\) in \(\mathfrak{g}\) and \(\lambda \in \mathbf{K}\) . A Lie \(p\) -algebra over \(\mathbf{K}\) is a set with the structure defined by giving it a Lie algebra structure over \(\mathbf{K}\) and a \(p\) -mapping. Every associative algebra over \(\mathbf{K}\) is a Lie \(p\) -algebra with \([x, y] = xy - yx\) and \(x^{[p]} = x^p\) (Exercise 19). A mapping \(u\) of a Lie \(p\) -algebra \(\mathfrak{g}\) into a Lie \(p\) -algebra \(\mathfrak{g}'\) is called a \(p\) -homomorphism if \(u\) is a Lie algebra homomorphism and \(u(x^{[p]}) = (u(x))^{[p]}\) . Show that in a Lie \(p\) -algebra \(\mathfrak{g}\) every \(p\) -mapping is of the form \(x \mapsto x^{[p]} + f(x)\) , where \(f\) is a mapping of \(\mathfrak{g}\) into its centre, which is semi-linear for the endomorphism \(\lambda \mapsto \lambda^p\) of \(\mathbf{K}\) . More generally, if \(u\) is a homomorphism of \(\mathfrak{g}\) onto a Lie \(p\) -algebra \(\mathfrak{g}', u(x^{[p]}) - (u(x))^{[p]}\) belongs to the centre of \(\mathfrak{g}'\) .

\begin{enumerate}
\def\labelenumi{\arabic{enumi}.}
\setcounter{enumi}{20}
\tightlist
\item
  \begin{enumerate}
  \def\labelenumii{(\alph{enumii})}
  \tightlist
  \item
    Let \(\mathbf{K}\) be a field of characteristic \(p > 0\) and \(\mathbf{L}\) a not necessarily associative algebra over \(\mathbf{K}\) . Show that the derivations of \(\mathbf{L}\) form a Lie \(p\) -algebra with the \(p\) -mapping \(\mathbf{D} \mapsto \mathbf{D}^p\) .
  \end{enumerate}
\end{enumerate}

\begin{enumerate}
\def\labelenumi{(\alph{enumi})}
\setcounter{enumi}{1}
\item
  Let \(\mathbf{L}\) be the algebra \(\mathrm{K}[\mathrm{X}] / (\mathrm{X}^p)\) and \(\mathfrak{w}\) the \(p\) -algebra of derivations of \(\mathbf{L}\) . Let \(\mathrm{D}_i\) ( \(0 \leqslant i \leqslant p - 1\) ) denote the derivation of \(\mathbf{L}\) such that \(\mathrm{D}_i(\mathbf{X}) = \mathbf{X}^i\) . Show that the \(\mathrm{D}_i\) form a basis of \(\mathfrak{w}\) over \(\mathbf{K}\) and that \([\mathrm{D}_i, \mathrm{D}_j] = (j - i)\mathrm{D}_{i + j - 1}\) if \(i + j \leqslant p\) , \([\mathrm{D}_i, \mathrm{D}_j] = 0\) if \(i + j > p\) , and \(\mathrm{D}_i^p = 0\) , except for \(i = 1\) , in which case \(\mathrm{D}_1^p = \mathrm{D}_1\) .
\item
  Show that, if \(p \geqslant 3\) , the Lie algebra \(\mathfrak{w}\) only admits the ideals \(\{0\}\) and \(\mathfrak{w}\) . (Using the multiplication table of the \(D_i\) , show that every non-zero ideal of \(\mathfrak{w}\) contains a non-zero multiple of \(D_{p-1}\) .)
\item
  Let \(V_{k}\) be the subspace of w with basis the \(D_{i}\) of index \(\geqslant k\) . Show that, if \(p \geqslant 5\) , \(V_k\) is identical with the set of \(Z \in w\) whose centralizer is of dimension \(\geqslant k + 1\) . Deduce that, if \(p \geqslant 5\) , the automorphism group of \(w\) is solvable.
\end{enumerate}

\begin{enumerate}
\def\labelenumi{\arabic{enumi}.}
\setcounter{enumi}{21}
\tightlist
\item
  Let g be a Lie p-group over a ring K of characteristic p (p prime). Let \(x \mapsto x^{[p]}\) denote the p-mapping. For every subset E of g, let \(E^{p}\) denote the submodule of g generated by the \(x^{p}\) , where \(x \in E\) . An ideal \(a \subset g\) is called a p-ideal if \(a^{p} \subset a\) .
\end{enumerate}

\begin{enumerate}
\def\labelenumi{(\alph{enumi})}
\item
  For an ideal \(\mathfrak{a} \subset \mathfrak{g}\) to be a \(p\) -ideal, it is necessary and sufficient that it be the kernel of a \(p\) -homomorphism (Exercise 20).
\item
  The sum of two \(p\) -ideals is a \(p\) -ideal. The sum of a \(p\) -ideal and a \(p\) -subalgebra is a \(p\) -subalgebra.
\item
  Let \(\mathfrak{h}\) be a Lie subalgebra of \(\mathfrak{g}\) . The smallest Lie \(p\) -subalgebra containing \(\mathfrak{h}\) is \(\mathfrak{h} + \mathfrak{h}^p + \mathfrak{h}^{p^2} + \cdots = \overline{\mathfrak{h}}\) ; if we write \(\mathfrak{h}_i = \mathfrak{h} + \mathfrak{h}^p + \cdots + \mathfrak{h}^{p^i}\) , \(\mathfrak{h}_i\) is an ideal of \(\overline{\mathfrak{h}}\) and \(\overline{\mathfrak{h}}/\mathfrak{h}\) is commutative. If \(\mathfrak{h}\) is an ideal of \(\mathfrak{g}\) , so are the \(\mathfrak{h}_i\) and \(\overline{\mathfrak{h}}\) and the latter is the smallest \(p\) -ideal containing \(\mathfrak{h}\) .
\item
  The normalizer and centralizer of an arbitrary submodule of g are p-subalgebras of g.
\end{enumerate}

\begin{enumerate}
\def\labelenumi{\arabic{enumi}.}
\setcounter{enumi}{22}
\tightlist
\item
  Let \(\mathfrak{g}\) be a finite-dimensional commutative Lie \(p\) -algebra over a perfect field \(\mathbf{K}\) of characteristic \(p > 0\) .
\end{enumerate}

\begin{enumerate}
\def\labelenumi{(\alph{enumi})}
\item
  Show that g is in a unique way the direct sum of two p-subalgebras h, t such that on h the p-mapping is bijective and on t it is nilpotent (consider the iterations of the p-mapping on g; cf.~Algebra, Chapter VIII, § 2, no. 2, Lemma 2 and Chapter VII, § 5, Exercise 20 (a)). h is called the p-core of g.
\item
  Suppose that K is algebraically closed. Show that there exists a basis \((e_{i})\) of \(\mathfrak{h}\) such that \(e_{i}^{p}=e_{i}\) for all i (consider a p-subalgebra of h generated by a single element and show that it contains some x such that \(x^{p}=x\neq0\) ; then argue by induction on the dimension of \(\mathfrak{h}\)). Deduce that the p-subalgebras of h are finite in number and are in one-to-one correspondence with the vector subspaces of the vector space over the prime field \(F_{p}\) generated by the \(e_{i}\) .
\item
  Show that there exists a basis \((f_{ij})\) of \(\mathfrak{t}\) ( \(1 \leqslant i \leqslant k\) , \(1 \leqslant j \leqslant s_i\) for all \(i\) ) where the sequence of \(s_i\) is decreasing and which is such that \(f_{1j}^p = 0\) for \(1 \leqslant j \leqslant s_1\) , \(f_{ij}^p = f_{i-1,j}\) for \(2 \leqslant i \leqslant k\) , \(1 \leqslant j \leqslant s_i\) (method of Algebra, Chapter VII, § 5, Exercise 20 (b)).
\end{enumerate}

¶ 24. Let K be a (commutative) field of (arbitrary) characteristic p and \(X_{i}, Y_{i}, Z_{i}\) 3n distinct indeterminates; to abbreviate we shall denote the systems \((\mathrm{X}_{i})_{1 \leq i \leq n}, (\mathrm{Y}_{i})_{1 \leq i \leq n}, (\mathrm{Z}_{i})_{1 \leq i \leq n}\) by x, y, z respectively. Let \(A_{x,y,z}\) be the algebra of formal power series in the 3n indeterminates \(X_{i}, Y_{i}, Z_{i}\) with coefficients in K; we shall denote by \(A_{x,y}\) (resp. \(A_{x}\) ) the subalgebra of \(A_{x,y,z}\) consisting of the power series not containing z (resp. y, z). An element of \(A_{x,y,z}\) (resp. \(A_{x,y}, A_{x}\) ) will be denoted by \(u(\mathbf{x}, \mathbf{y}, \mathbf{z})\) (resp. \(u(\mathbf{x}, \mathbf{y}), u(\mathbf{x})\) );

 a system \((u_{i}(\mathbf{x},\mathbf{y},\mathbf{z}))_{1\leqslant i\leqslant n}\) of n elements of \(A_{x,y,z}\) will be denoted by \(u(\mathbf{x},\mathbf{y},\mathbf{z})\) ; analogous notation will be used for formal power series containing only one or two of the three systems x, y, z. If \(u(\mathbf{x},\mathbf{y},\mathbf{z})\in A_{x,y,z}\) and \(\mathbf{f}=(f_{i}),\mathbf{g}=(g_{i}),\mathbf{h}=(h_{i})\) are three systems of n elements of \(A_{x,y,z}\) which are formal power series with no constant term, we denote by \(u(\mathbf{f}(\mathbf{x},\mathbf{y},\mathbf{z}),\mathbf{g}(\mathbf{x},\mathbf{y},\mathbf{z}),\mathbf{h}(\mathbf{x},\mathbf{y},\mathbf{z}))\) the formal power series obtained by substituting \(f_{i}\) for \(X_{i}, g_{i}\) for \(Y_{i}, h_{i}\) for \(Z_{i}\) ( \(1\leqslant i\leqslant n\) ) in u. For every system \(\alpha=(\alpha_{1},\ldots,\alpha_{n})\) of n integers \(\geqslant0\) , let \(x^{\alpha}\) denote the monomial \(X_{1}^{\alpha_{1}}\ldots X_{n}^{\alpha_{n}}\) ; \(y^{\alpha}\) and \(z^{\alpha}\) are defined similarly. Let \(\varepsilon_{i}\) denote the system \(\alpha\) for which \(\alpha_{j}=0\) if \(j\neq i, \alpha_{i}=1\) . We denote by e the system \((0,\ldots,0)\) of n elements of \(A_{x,y,z}\) .

\begin{enumerate}
\def\labelenumi{(\alph{enumi})}
\item
  A formal group law over K (or, by an abuse of language, a formal group over K) of dimension n is a system \(G = \mathbf{f}(\mathbf{x}, \mathbf{y})\) of n elements of \(A_{x,y}\) with the following properties: (1) \(\mathbf{f}(\mathbf{x}, \mathbf{f}(\mathbf{y}, \mathbf{z})) = \mathbf{f}(\mathbf{f}(\mathbf{x}, \mathbf{y}), \mathbf{z})\) ; (2) \(\mathbf{f}(\mathbf{e}, \mathbf{y}) = \mathbf{y}\) , \(\mathbf{f}(\mathbf{x}, \mathbf{e}) = \mathbf{x}\) . Show that necessarily \(f_{i}(\mathbf{x}, \mathbf{y}) = \mathrm{X}_{i} + \mathrm{Y}_{i} + g_{i}(\mathbf{x}, \mathbf{y})\) , where \(g_{i}\) contains only monomials of degree \(\geqslant 2\) , each of which contains at least one \(X_{j}\) and at least one \(Y_{j}\) . Show that there exists one and only one system \(\mathbf{h}(\mathbf{x})\) of n elements of \(A_{x}\) such that \(\mathbf{f}(\mathbf{x}, \mathbf{h}(\mathbf{x})) = \mathbf{f}(\mathbf{h}(\mathbf{x}), \mathbf{x}) = \mathbf{e}\) (Algebra, Chapter IV, § 5, no. 9, Proposition 10). G is called commutative if \(\mathbf{f}(\mathbf{y}, \mathbf{x}) = \mathbf{f}(\mathbf{x}, \mathbf{y})\) .
\item
  For all \(u \in \mathbf{A}_{\mathbf{x}}\) , let \(L_{\mathbf{y}}u\) denote the element \(u(\mathbf{f}(\mathbf{y},\mathbf{x}))\) of \(\mathbf{A}_{\mathbf{x},\mathbf{y}}\) . A derivation D of \(\mathbf{A}_{\mathbf{x}}\) can be canonically extended to a derivation of \(\mathbf{A}_{\mathbf{x},\mathbf{y}}\) (also denoted D by an abuse of notation) by the condition that \(\mathrm{D}(\mathrm{Y}_i) = 0\) for all \(i\) (Algebra, Chapter IV, § 5, no. 8, Proposition 6). D is said to be left invariant for the formal group G in question if \(L_{\mathbf{y}}D = DL_{\mathbf{y}}\) . We denote by \(\mathrm{D}^{(\mathrm{e})}\) the linear mapping of \(\mathbf{A}_{\mathbf{x},\mathbf{y}}\) into \(\mathbf{A}_{\mathbf{y}}\) defined by \(\mathrm{D}^{(\mathrm{e})}(u) = (\mathrm{D}u)(\mathbf{e},\mathbf{y})\) ; it maps \(\mathbf{A}_{\mathbf{x}}\) into K and is determined by its restriction \(\mathbf{A}_{\mathbf{x}}\) . Show that, for D to be left invariant, it is necessary and sufficient that for all \(v \in \mathbf{A}_{\mathbf{x}}\) , \(\mathrm{D}^{(\mathrm{e})}(L_{\mathbf{y}}v) = (\mathrm{D}v)(\mathbf{y})\) .
\end{enumerate}

Let \(D_{i}\) be the derivation \(\partial/\partial X_{i}\) of \(A_{x}\) ( \(1 \leqslant i \leqslant n\) ); there exists one and only one left invariant derivation \(T_{i}\) such that \(T_{i}^{(e)} = D_{i}^{(e)}\) . Show that the \(T_{i}\) are linearly independent over K and that every left invariant derivation is a linear combination of the \(T_{i}\) with coefficients in K. Deduce that, for the bracket {[}D, D'{]} and the p-mapping \(D \mapsto D^{p}\) (when p \textgreater{} 0), the set g of left invariant derivations is a Lie algebra (a Lie p-algebra if p \textgreater{} 0) called the Lie algebra of the formal Lie group G.

\begin{enumerate}
\def\labelenumi{(\alph{enumi})}
\setcounter{enumi}{2}
\tightlist
\item
  Show that for every formal power series \(u \in \mathbf{A}_{\mathbf{x}}\) :
\end{enumerate}

\[
u (\mathbf {f} (\mathbf {x}, \mathbf {y})) = u (0) + \sum_ {i = 1} ^ {n} Y _ {i} T _ {i} u + o _ {2} (\mathbf {x}, \mathbf {y})\tag{1}
\]

where in the series \(o_{2}\) all the monomials are of degree \(\geqslant2\) in the \(Y_{j}\) . Deduce that

\[
\begin{array}{r l} {u (\mathbf {f} (\mathbf {x}, \mathbf {f} (\mathbf {y}, \mathbf {z}))) = u (0) + \sum_ {i = 1} ^ {n} (\mathrm{Y} _ {i} + \mathrm{Z} _ {i}) \mathrm{T} _ {i} u} & \\ {+ \sum_ {i, j} \mathrm{Y} _ {i} \mathrm{Z} _ {j} ((\mathrm{T} _ {i} \mathrm{T} _ {j}) (u)) + o _ {3} (\mathbf {x}, \mathbf {y}, \mathbf {z})} & \end{array}\tag{2}
\]

where in the series \(o_{3}\) all the monomials are of degree \(\geqslant3\) in the set of \(Y_{i}\) and \(Z_{i}\) . Deduce that:

\[
u (\mathbf {f} (\mathbf {x}, \mathbf {y})) - u (\mathbf {f} (\mathbf {y}, \mathbf {x})) = \sum_ {i <   j} (X _ {i} Y _ {j} - X _ {j} Y _ {i}) ([ T _ {i}, T _ {j} ] ^ {(e)} (u)) + o _ {3} ^ {\prime} (\mathbf {x}, \mathbf {y})\tag{3}
\]

where all the terms of \(o_3'\) are of total degree \(\geqslant 3\) . Show that, if \(G\) is commutative, \(g\) is commutative.

\begin{enumerate}
\def\labelenumi{(\alph{enumi})}
\setcounter{enumi}{3}
\tightlist
\item
  Let G' be another formal group of dimension m, whose group law is f'. A formal homomorphism of G into G' is a system \(\mathbf{F} = (\mathrm{F}_{j}(\mathbf{x}))_{1 \leqslant j \leqslant m}\) of elements of \(\mathbf{A}_{\mathbf{x}}\) such that \(\mathbf{f}'(\mathbf{F}(\mathbf{x}), \mathbf{F}(\mathbf{y})) = \mathbf{F}(\mathbf{f}(\mathbf{x}, \mathbf{y}))\) . For every left invariant derivation \(D \in g\) , show that \(v \mapsto D(v \circ F)\) is a left invariant derivation belonging to the Lie algebra \(g'\) of G'. If it is denoted by \(\mathbf{F}^*(D)\) , \(\mathbf{F}^*\) is a homomorphism (a p-homomorphism for p \textgreater{} 0) of g into \(g'\) . If \((\mathrm{T}_{1})_{1 \leqslant i \leqslant n}, (\mathrm{T}'_{j})_{1 \leqslant j \leqslant m}\) are the bases of g and \(g'\) such that \(T_{i}^{(e)} = D_{i}^{(e)}\) and \(T'_{j}{}^{(e)} = D'_{j}{}^{(e)}\) , show that the matrix of \(F^*\) relative to these bases is the matrix \((\partial F_{j}/\partial X_{i})_{0}\) (constant term of \(\partial F_{j}/\partial X_{i}\) , where i is the index of the rows and j that of the columns).
\end{enumerate}

\begin{enumerate}
\def\labelenumi{\arabic{enumi}.}
\setcounter{enumi}{24}
\tightlist
\item
  The formal linear group in \(n\) variables over the field \(\mathbf{K}\) , denoted by \(\mathbf{GL}(n, \mathbf{K})\) (when no confusion can arise), is the formal group of dimension \(n^2\) whose group law is defined by \(f_{ij}(\mathbf{u}, \mathbf{v}) = -\delta_{ij} + \sum_{k=1}^{n} (\delta_{ik} + \mathrm{U}_{ik})(\delta_{kj} + \mathrm{V}_{kj})\) ( \(\delta_{ij}\) the Kronecker index; the \(3n^2\) indeterminates of the general definition of Exercise 24 are here denoted by \(\mathrm{U}_{ij}, \mathrm{V}_{ij}, \mathrm{W}_{ij}\) with 2 indices varying from 1 to \(n\) ). Show that, if \(\mathrm{D}_{ij} = \partial/\partial \mathrm{U}_{ij}\) , the left invariant derivations \(\mathrm{X}_{ij}\) such that \(\mathrm{X}_{ij}^{(e)} = \mathrm{D}_{ij}^{(e)}\) are given by
\end{enumerate}

\[
\mathrm{X} _ {i j} = (1 + \mathrm{U} _ {i i}) \mathrm{D} _ {i j} + \sum_ {k \neq i} \mathrm{U} _ {k i} \mathrm{D} _ {k j}.
\]

The Lie algebra of \(\mathbf{GL}(n,\mathbf{K})\) is identified with \(\mathfrak{gl}(n,\mathbf{K})\) by identifying \(X_{ij}\) with the element \(E_{ij}\) of the canonical basis (use formula (3) of Exercise 24). If K is of characteristic p \textgreater{} 0, \(X_{ii}^{p} = X_{ii}\) and \(X_{ij}^{p} = 0\) for \(i \neq j\) .

¶ 26. (a) Let K be a field of characteristic \(\neq 2\) , E an \(n\) -dimensional vector space over K and \(\Phi\) a non-degenerate symmetric bilinear form on E. Suppose that E has an orthogonal basis with respect to \(\Phi\) and let \(R\) denote the (diagonal) matrix of \(\Phi\) with respect to this basis. We know (Algebra, Chapter IX, § 4, Exercise 11) that every matrix (relative to the basis in question ) \(U\) of the orthogonal group \(\mathbf{G}(\Phi)\) such that \(\det(I + U) \neq 0\) can be written

\[
U = (I - R ^ {- 1} S) ^ {- 1} (I + R ^ {- 1} S),
\]

  where \(S = R(U - I)(U + I)^{-1}\) is an alternating matrix such that \(\det(I - R^{-1}S) \neq 0\) ; conversely, for every matrix \(S\) satisfying these conditions, \(U = (I - R^{-1}S)^{-1}(I + R^{-1}S)\) is a matrix of \(\mathbf{G}(\Phi)\) such that \(\det(I + U) \neq 0\) . Let \(S_{ij}, S_{ij}', S_{ij}''\) be \(3n(n - 1)/2\) indeterminates ( \(1 \leqslant i < j \leqslant n\) )

and let L be a field containing the ring of formal power series \(K[[S_{ij}, S'_{ij}, S''_{ij}]]\) . If we denote by \(S, S', S''\) the matrices

\[
\sum_ {i <   j} \mathrm{S} _ {i j} (E _ {i j} - E _ {j i}), \qquad \sum_ {i <   j} \mathrm{S} _ {i j} ^ {\prime} (E _ {i j} - E _ {j i}), \qquad \sum_ {i <   j} \mathrm{S} _ {i j} ^ {\prime \prime} (E _ {i j} - E _ {j i}),
\]

then \(\operatorname{det}(I - R^{-1}S) \neq 0\) and the analogues for \(S'\) and \(S''\) hold; if we write

\[
U = (I - R ^ {- 1} S) ^ {- 1} (I + R ^ {- 1} S), \quad U ^ {\prime} = (I - R ^ {- 1} S ^ {\prime}) ^ {- 1} (I + R ^ {- 1} S ^ {\prime}),
\]

then also \(\operatorname{det}(I + UU') \neq 0\) . We write

\[
F (S, S ^ {\prime}) = R \left(U U ^ {\prime} - I\right) \left(U U ^ {\prime} + I\right) ^ {- 1} = \left(f _ {i j} \left(S, S ^ {\prime}\right)\right);
\]

the \(f_{ij}(S, S')\) belong to the ring of formal power series \(K[[S_{ij}, S_{ij}']]\) and, by considering \(F(S, S')\) as a formal power series with coefficients in the algebra of matrices of order n over K, we can write:

\[
F (S, S ^ {\prime}) = S + S ^ {\prime} + o _ {2} (S, S ^ {\prime})
\]

where in the elements of the matrix \(o_2(S, S')\) the terms are all of degree \(\geqslant 1\) in the \(S_{ij}\) and of degree \(\geqslant 1\) in the \(S_{ij}'\) ; then similarly:

\[
U = I + 2 R ^ {- 1} S + o _ {2} ^ {\prime} (S)
\]

the elements of \(o_2'(S)\) being formal power series of order \(\geqslant 2\) . Show that \(F(S, S')\) is a formal group law over \(K\) of dimension \(n(n - 1)/2\) ; this formal group is called the formal orthogonal group (corresponding to \(\Phi\) ) and is denoted by \(G(\Phi)\) by an abuse of notation.

\begin{enumerate}
\def\labelenumi{(\alph{enumi})}
\setcounter{enumi}{1}
\tightlist
\item
  If we write \(M(S) = (m_{ij}(S)) = (I - R^{-1}S)^{-1}(I + R^{-1}S)\) , \(M\) is a formal homomorphism of \(\mathbf{G}(\Phi)\) into \(\mathbf{GL}(n, \mathbf{K})\) . Let \(\mathfrak{o}(\Phi)\) be the Lie algebra of \(\mathbf{G}(\Phi)\) and let \((\mathrm{T}_{ij})_{i < j}\) denote the basis of this algebra such that \(\mathrm{T}_{ij}^{(e)} = \mathrm{D}_{ij}^{(e)}\) (cf.~Exercise 24 (d)); if an element \(D = \sum_{i < j} c_{ij} T_{ij}\) of \(\mathfrak{o}(\Phi)\) is identified with the matrix \(D = \sum_{i < j} c_{ij}(E_{ij} - E_{ji})\) , show that (in the notation of Exercise 24 (d) with the identification made in Exercise 25 of the Lie algebra of \(\mathbf{GL}(n, \mathbf{K})\) with \(\mathfrak{gl}(n, \mathbf{K})\) ):
\end{enumerate}

\[
M ^ {*} (D) = 2 R ^ {- 1} D.
\]

  Deduce that \(M^*\) is an isomorphism of \(\mathfrak{o}(\Phi)\) onto the subalgebra of \(\mathfrak{gl}(n, \mathbf{K})\) consisting of the matrices \(X\) such that \({}^{t}XR + RX = 0\) (for every ordered pair \((a, b)\) of elements of \(E\) , identified with matrices with one column, consider \(\Phi(a + M(S).a, b + M(S).b) = {}^{t}a.(I + {}^{t}M(S))R(I + M(S)).b\) as a formal power series reduced to its constant term and use the fact that the subalgebra of \(X \in \mathfrak{gl}(n, \mathbf{K})\) such that \({}^{t}XR + RX = 0\) is of dimension \(n(n - 1)/2\) ).

\begin{enumerate}
\def\labelenumi{(\alph{enumi})}
\setcounter{enumi}{2}
\tightlist
\item
  Define similarly the formal symplectic group in 2n variables over an arbitrary
\end{enumerate}

 (commutative) field and show that its Lie algebra is identified with the subalgebra of \(\mathfrak{gl}(2n,\mathbf{K})\) consisting of the matrices X such that

\[
{ } ^ { t } X A + A X = 0 , \quad \text {   where   } \quad A = \left( \begin{array} { c c } 0 & I _ { n } \\ - I _ { n } & 0 \end{array} \right) .
\]

\begin{enumerate}
\def\labelenumi{\arabic{enumi}.}
\setcounter{enumi}{26}
\tightlist
\item
  Let K be a field of characteristic p \textgreater{} 0 and G the 2-dimensional formal group defined by \(f_{1}(\mathbf{x}, \mathbf{y}) = \mathrm{X}_{1} + \mathrm{Y}_{1} + \mathrm{X}_{1}\mathrm{Y}_{1}, f_{2}(\mathbf{x}, \mathbf{y}) = \mathrm{X}_{2} + \mathrm{Y}_{2}(1 + \mathrm{X}_{1}^{p})\) . Show that G is not commutative but that its Lie algebra is commutative.
\end{enumerate}

\exercisesfor{2}

\begin{enumerate}
\def\labelenumi{\arabic{enumi}.}
\item
  We adopt the notation T, J, \(T_{n}\) of Definition 1 and no. 6. Let t be a homogeneous tensor of T of order p and \(\sigma\) a permutation of \(\{1, 2, \ldots, p\}\) . Then \(t - \sigma t \in J + T_{p-1}\) . (Reduce it to the case where \(\sigma\) is a transposition of two consecutive integers.)
\item
  Suppose that \(\mathbf{K}\) is a field. Let \(\mathfrak{g}\) be a Lie \(\mathbf{K}\) -algebra and \(\mathbf{U}\) its enveloping algebra.
\end{enumerate}

\begin{enumerate}
\def\labelenumi{(\alph{enumi})}
\item
  Let \(u, v\) be in U. If \(u\) is of filtration \(> n\) , and \(v\) of filtration \(> p\) , then \(uv\) is of filtration \(> n + p\) (use Theorem 1).
\item
  Deduce from (a) that the only invertible elements of U are the scalars.
\item
  Deduce from (b) that the Jacobson radical of U is zero.
\end{enumerate}

\begin{enumerate}
\def\labelenumi{\arabic{enumi}.}
\setcounter{enumi}{2}
\item
  Suppose that K is a field. Let g be a Lie K-algebra and U its enveloping algebra. The left regular representation of U corresponds to a representation \(\rho\) of g on the space U. Show that the set \(U_{+}\) of elements of U with no constant term is stable under \(\rho\) and that \(U_{+}\) has no supplement in U stable under \(\rho\) . In particular, \(\rho\) is not semi-simple. (This will be considered again in Theorem 2 of § 6.)
\item
  Suppose that K is a field.
\end{enumerate}

\begin{enumerate}
\def\labelenumi{(\alph{enumi})}
\item
  Verify that there exists a 3-dimensional Lie algebra g with basis \((x, y, z)\) such that \([x, y] = z\) , \([x, z] = [y, z] = 0\) . Let U be the enveloping algebra of g. Show that the centre of U is the subalgebra generated by 1 and z. (Consider, for all \(\xi \in g\) , the derivation of the symmetric algebra S of g which extends ad, \(\xi\) and look for the elements of S annihilated by these derivations; then apply the final remark of no. 8.)
\item
  Verify that there exists a 3-dimensional Lie algebra g with basis \((x, y, z)\) such that \([x, y] = y\) , \([x, z] = z\) , \([y, z] = 0\) . Show that the centre of the enveloping algebra of g reduces to the scalars. (Same method.)
\end{enumerate}

\begin{enumerate}
\def\labelenumi{\arabic{enumi}.}
\setcounter{enumi}{4}
\item
  Suppose that \(\mathbf{K}\) is a field of characteristic \(p > 0\) ( \(p\) prime). Let \(\mathfrak{g}\) be a Lie algebra over \(\mathbf{K}\) with a basis, canonically identified with a submodule of its enveloping algebra U. For an endomorphism \(\phi\) of the K-module \(\mathfrak{g}\) to be a \(p\) -mapping, it is necessary and sufficient that \(x \mapsto x^p - \phi(x)\) be a semi-linear mapping (for \(\lambda \mapsto \lambda^p)\) of \(\mathfrak{g}\) into the centre of U. Deduce that if \((b_{\lambda})\) is a basis of \(\mathfrak{g}\) , for there to exist a \(p\) -mapping on \(\mathfrak{g}\) , it is necessary and sufficient that, for all \(\lambda\) , there exist \(c_{\lambda} \in \mathfrak{g}\) such that \((\operatorname{ad} b_{\lambda})^p = \operatorname{ad} c_{\lambda}\) ; if this is so, there then exists one and only one \(p\) -mapping \(x \mapsto x^{[p]}\) such that \(b_{\lambda}^{[p]} = c_{\lambda}\) for all \(\lambda\) .
\item
  Let g be a Lie p-algebra over a ring K such that \(pK = \{0\}\) (p prime), U its enveloping algebra and \(\sigma\) the canonical mapping of g into U. Let J be the two-sided ideal of U generated by the elements \((\sigma(x))^{p} - \sigma(x^{[p]})\) where x runs through g. The associative algebra \(\tilde{U} = U/J\) is called the restricted enveloping algebra of g. The mapping \(\sigma\) defines when passing to the quotient a mapping \(\tilde{\sigma}\) (called canonical) of g into \(\tilde{U}\) , which is a p-homomorphism (when \(\tilde{U}\) is considered as a Lie p-algebra).
\end{enumerate}

\begin{enumerate}
\def\labelenumi{(\alph{enumi})}
\item
  The algebra \(\tilde{\mathbf{U}}\) and the mapping \(\tilde{\sigma}\) are solutions of a universal mapping problem: for every \(p\) -homomorphism \(f\) of \(\mathfrak{g}\) into an associative algebra \(\mathbf{B}\) over \(\mathbf{K}\) (considered as a Lie \(p\) -algebra), there exists one and only one K-homomorphism \(f'\) of \(\tilde{\mathbf{U}}\) into \(\mathbf{B}\) (with their associative algebra structures) such that \(f = f' \circ \tilde{\sigma}\) .
\item
  Show that, if \((b_{\lambda})_{\lambda \in \Lambda}\) is a basis of \(\mathfrak{g}\) (where \(\Lambda\) is totally ordered), \(\tilde{\sigma}\) is injective and that, if \(x\) and \(\tilde{\sigma}(x)\) are identified for \(x\in \mathfrak{g}\) , the elements \(\Pi_{\lambda}b_{\lambda}^{\nu_{\lambda}}\) (where the \(\lambda\) are in increasing order, the \(\nu_{\lambda}\) are zero except for a finite number and \(0\leqslant \nu_{\lambda} < p\) for all \(\lambda\) ) form a basis of \(\tilde{\mathbf{U}}\) . (Canonically identifying \(\mathfrak{g}\) with a submodule of \(\mathrm{U}\) , we write \(\phi (x) = x^{p} - x^{[p]}\) for all \(x\in g\) . For every composite index \(\alpha = (\alpha_{\lambda})\in \mathbf{N}^{(\Lambda)}\) , let \(\alpha_{\lambda} = \beta_{\lambda} + p\gamma_{\lambda}\) with \(0\leqslant \beta_{\lambda} < p\) and let
\end{enumerate}

\[
\mathrm{T} _ {\alpha} = (\Pi_ {\lambda} b _ {\lambda} ^ {\beta_ {\lambda}}) (\Pi_ {\lambda} (\phi (b _ {\lambda})) ^ {\gamma_ {\lambda}}).
\]

Show that the \(T_{\alpha}\) form a basis of U and the \(T_{\alpha}\) such that \(\gamma = (\gamma_{\lambda}) \neq 0\) form a basis of J; observe that the \(\phi(b_{\lambda})\) belong to the centre of U.)

\begin{enumerate}
\def\labelenumi{\arabic{enumi}.}
\setcounter{enumi}{6}
\tightlist
\item
  Let g be a Lie p-algebra over a ring K such that \(pK = \{0\}\) (p prime). A derivation D of g is called a p-derivation if \(\mathrm{D}(x^{p}) = (\mathrm{ad} x)^{p-1}\) . Dx for all \(x \in g\) . Every inner derivation is a p-derivation.
\end{enumerate}

\begin{enumerate}
\def\labelenumi{(\alph{enumi})}
\item
  If \(\mathbf{L}\) is an associative K-algebra, every derivation of \(\mathbf{L}\) is a \(p\) -derivation when \(\mathbf{L}\) is considered as a Lie \(p\) -algebra. (Use formula (2) of Exercise 19 of § 1.)
\item
  Suppose that g has a basis. For a derivation of g to be a p-derivation, it is necessary and sufficient that it can be extended to a derivation of the restricted enveloping algebra of g. Deduce that the p-derivations of g form a Lie p-subalgebra of the p-algebra of derivations of g.
\item
  If D is a \(p\) -derivation of \(\mathfrak{g}\) , the kernel of D is a \(p\) -subalgebra of \(\mathfrak{g}\) .
\item
  For every derivation D of g, \(\mathrm{D}(x^{p}) - (\mathrm{ad} x)^{p-1}\) . \(\bar{D}x\) belongs to the centre of g for all \(x \in g\) (use formula (2) of Exercise 19 of §1, applied to \(\mathrm{L} = \mathcal{L}(\mathfrak{g})\) ).
\end{enumerate}

\begin{enumerate}
\def\labelenumi{\arabic{enumi}.}
\setcounter{enumi}{7}
\item
  Show that Theorem 1 remains valid if the module \(\mathfrak{g}\) is a direct sum of monogenous modules. (Replace the module P in the proof by the symmetric algebra of g.)
\item
  Let \(k\) be the field with two elements, \(\mathbf{V}\) the vector space \(k^3\) , \((x_1, x_2, x_3)\) its canonical basis and \(\mathbf{K}\) the exterior algebra of \(\mathbf{V}\) , which is an 8-dimensional commutative algebra over \(k\) . Let \(\mathfrak{g}\) be the Lie K-algebra admitting the basis \((e_1, e_2, e_3, e_{12}, e_{13}, e_{23})\) such that
\end{enumerate}

\[
\left[ e _ {1}, e _ {2} \right] = \left[ e _ {2}, e _ {1} \right] = e _ {1 2} \quad \left[ e _ {1}, e _ {3} \right] = \left[ e _ {3}, e _ {1} \right] = e _ {1 3}, \quad \left[ e _ {2}, e _ {3} \right] = \left[ e _ {3}, e _ {2} \right] = e _ {2 3}
\]

and the other brackets are zero. Let \(\mathfrak{h}\) be the ideal of \(\mathfrak{g}\) generated by \(u = x_{1}e_{1} + x_{2}e_{2} + x_{3}e_{3}\) . As a module, \(\mathfrak{h}\) is generated by \(u\) , \([u, e_1] = x_2e_{12} + x_3e_{13}\) , \([u, e_2] = x_1e_{12} + x_3e_{23}\) , \([u, e_3] = x_1e_{13} + x_2e_{23}\) . Let

\[
v = x _ {1} x _ {2} e _ {1 2} + x _ {1} x _ {3} e _ {1 3} + x _ {2} x _ {3} e _ {2 3}.
\]

\[
\phi \left(e _ {1}\right) = \phi \left(e _ {2}\right) = \phi \left(e _ {3}\right) = 0, \phi \left(e _ {1 2}\right) = x _ {3}, \phi \left(e _ {1 3}\right) = x _ {2}, \phi \left(e _ {2 3}\right) = x _ {1}.
\]

\begin{enumerate}
\def\labelenumi{(\alph{enumi})}
\setcounter{enumi}{1}
\item
  Let \(f\) be a linear mapping of \(\mathfrak{g}\) into an associative K-algebra such that \(f([x,y]) = f(x)f(y) - f(y)f(x)\) for all \(x, y\) in \(\mathfrak{g}\) . Show that \(f(v) = f(u)^2\) .
\item
  Deduce from (a) and (b) that the canonical mapping of the Lie algebra \(\mathfrak{g} / \mathfrak{h}\) into its enveloping algebra is not injective.
\end{enumerate}

\begin{enumerate}
\def\labelenumi{\arabic{enumi}.}
\setcounter{enumi}{9}
\tightlist
\item
  Let \(\mathfrak{g}\) be a finite-dimensional Lie algebra over a field, U its enveloping algebra and \(\mathbf{U}_n\) the set of elements of U of filtration \(\leqslant n\) .
\end{enumerate}

\begin{enumerate}
\def\labelenumi{(\alph{enumi})}
\item
  Let \(x \in \mathrm{U}_n, y \in \mathrm{U}_n\) be two non-zero elements. Show that there exists \(u \in \mathrm{U}, v \in \mathrm{U}\) such that \(ux = vy\) . (Compare \(\dim(\mathrm{U}_p x) = \dim \mathrm{U}_p\) , \(\dim(\mathrm{U}_p y) = \dim \mathrm{U}_p\) and \(\dim \mathrm{U}_{p+n}\) . Deduce that \(\mathrm{U}_p x \cap \mathrm{U}_p y \neq \{0\}\) for \(p\) sufficiently large.)
\item
  Show that U admits a field of left quotients (Algebra, Chapter I, § 9, Exercise 8) which is at the same time a field of right quotients.
\end{enumerate}

\exercisesfor{3}

\begin{enumerate}
\def\labelenumi{\arabic{enumi}.}
\item
  Let g be a Lie algebra, \(\rho\) a representation of g on a K-module M and \(\sigma\) the associated representation of g on the tensor algebra of M. Show that the submodule of symmetric tensors and the submodule of skew-symmetric tensors are stable under \(\sigma\) .
\item
  Let g be a Lie algebra, \(\rho\) a representation of g on a K-module M and \(\sigma\) the associated representation of g on Q = L(M, M; M). For \(f \in Q\) to be invariant under \(\sigma\) , it is necessary and sufficient that the \(\rho(x)\) be derivations of M with the multiplication defined by f.
\item
  Let V be a finite-dimensional vector space over a perfect field K.
\end{enumerate}

\begin{enumerate}
\def\labelenumi{(\alph{enumi})}
\tightlist
\item
  The identity representation of \(\mathfrak{gl}(\mathrm{V})\) defines canonically representations of \(\mathfrak{gl}(\mathrm{V})\) on \(\bigotimes_{q}^{p}\mathrm{V}\) and \(\bigotimes_{q}^{q}\mathrm{V}^{*}\) and hence a representation \(u\mapsto u_q^p\) of \(\mathfrak{gl}(\mathrm{V})\) on
\end{enumerate}

\(V_{q}^{p} = \left( \bigotimes_{i=1}^{p} V \right) \otimes \left( \bigotimes_{i=1}^{q} V^{*} \right)\) . Show that \(u_{q}^{p}\) is semi-simple if \(u\) is. (Reduce it, by extending the base field, to the case where \(K\) is algebraically closed.) Show that \(u_{q}^{p}\) is nilpotent if \(u\) is. Deduce that, if \(s\) and \(n\) are the semi-simple and nilpotent components of \(u, u_{q}^{p}\) and \(n_{q}^{p}\) are the semi-simple and nilpotent components of \(u_{q}^{p}\) .

\begin{enumerate}
\def\labelenumi{(\alph{enumi})}
\setcounter{enumi}{1}
\item
  Deduce from (a) that, if an element of \(V_q^p\) is annihilated by \(u_q^p\) , it is annihilated by \(s_q^p\) and \(n_q^p\) .
\item
  Deduce from (b) and Exercise 2 that, if V has a not necessarily associative algebra structure and u is a derivation of V, then s and n are derivations of V.
\end{enumerate}

\begin{enumerate}
\def\labelenumi{\arabic{enumi}.}
\setcounter{enumi}{3}
\tightlist
\item
  Suppose that K is a field. Let g be a Lie K-algebra.
\end{enumerate}

\begin{enumerate}
\def\labelenumi{(\alph{enumi})}
\item
  Let M and N be g-modules, where N is finite-dimensional over K. Let \((f_{j})_{1 \leqslant j \leqslant n}\) be a basis of N over K. If there are elements \(e_{1}, \ldots, e_{n}\) of M, not all zero such that \(\sum_{i=1}^{n} e_i \otimes f_i\) is invariant in \(\mathbf{M} \otimes \mathbf{N}\) , then the subspace \(\mathbf{M}_1\) of \(\mathbf{M}\) generated by the \(e_i\) is stable under the \(x_{\mathbf{M}}\) ; if further \(\mathbf{N}\) is simple the representation of \(\mathfrak{g}\) on \(\mathbf{M}_1\) is dual to the representation of \(\mathfrak{g}\) on \(\mathbf{N}\) .
\item
  Let \(M_{1}\) , \(M_{2}\) , N, \(N^{*}\) be g-modules of finite dimension over K. Suppose that \(M_{1}\) and \(M_{2}\) are simple and that the representations of g on N and \(N^{*}\) are dual. For the g-module \(M_{1}\) to be isomorphic to a sub-g-module of \(M_{2} \otimes N\) , it is necessary and sufficient that \(M_{2}\) be isomorphic to a sub-g-module of \(M_{1} \otimes N^{*}\) . (Use Proposition 4; consider the representation of g on \(M_{1}^{*} \otimes M_{2} \otimes N\) and apply (a) to the dual representation of the latter.)
\end{enumerate}

\begin{enumerate}
\def\labelenumi{\arabic{enumi}.}
\setcounter{enumi}{4}
\tightlist
\item
  Suppose that K is a field of characteristic 0. Let V be a finite-dimensional vector K-space, \(\mathfrak{g} = \mathfrak{sl}(V)\) , U the enveloping algebra of \(\mathfrak{g}\) , \(U_n\) the set of elements of filtration \(\leqslant n\) in U, \(U^n\) the set of images in U of homogeneous symmetric tensors of order \(n\) over \(\mathfrak{g}\) and \((\alpha_1, \ldots, \alpha_m)\) a basis of the vector space g. Let \(W_{s} = \bigotimes^{s} V\) . The identity representation of g defines a representation of g on \(W_{s}\) which extends to a homomorphism \(\pi_s\) of U into the algebra \(M_s = \mathcal{L}(W_s)\) .
\end{enumerate}

\begin{enumerate}
\def\labelenumi{(\alph{enumi})}
\tightlist
\item
  Let \(M_{s}^{\prime}\) be the subspace of \(M_{s}\) generated by the elements of the form \(\beta_{1} \otimes \beta_{2} \otimes \cdots \otimes \beta_{s}\) , where \(\beta_{i} \in \mathcal{L}(V)\) and \(\beta_{i} = 1\) for at least one i. Let z be an element of \(U_{s}\) and \(z_{0} = \sum_{1 \leqslant i_{1}, \ldots, i_{s} \leqslant m} a_{i_{1} \ldots i_{s}} \alpha_{i_{1}} \ldots \alpha_{i_{s}}\) its component in \(U^{s}\) , where the \(a_{i_{1} \ldots i_{s}}\) are symmetric with respect to permutation of indices. Show that
\end{enumerate}

\[
\pi_ {s} (z) \equiv s! \sum_ {1 \leqslant i _ {1}, \dots , i _ {s} \leqslant m} a _ {i _ {1} \dots i _ {s}} \alpha_ {i _ {1}} \otimes \dots \otimes \alpha_ {i _ {s}} \pmod {\mathrm{M} _ {s} ^ {\prime}}.
\]

\begin{enumerate}
\def\labelenumi{(\alph{enumi})}
\setcounter{enumi}{1}
\tightlist
\item
  Deduce from (a) that for all \(z \in \mathbf{U}\) there exists a finite-dimensional representation \(\pi\) of \(\mathbf{U}\) such that \(\pi(z) \neq 0\) . (For a sequel to this exercise, cf.~§ 7, Exercise 3.)
\end{enumerate}

\begin{enumerate}
\def\labelenumi{\arabic{enumi}.}
\setcounter{enumi}{5}
\tightlist
\item
  Suppose that \(\mathbf{K}\) is a field. Let \(\mathfrak{g}\) be a Lie K-algebra, \(x \mapsto x_{\mathbf{M}}\) a finite-dimensional representation of \(\mathfrak{g}\) , \((e_1, \ldots, e_n)\) a basis of \(\mathbf{M}\) and \((e_1^*, \ldots, e_n^*)\) the dual basis.
\end{enumerate}

\begin{enumerate}
\def\labelenumi{(\alph{enumi})}
\tightlist
\item
  If \(f\) is a linear form on \(\bigotimes^{r} \mathbf{M}\) , then:
\end{enumerate}

\[
f = \sum_ {1 \leqslant i _ {1} \dots , i _ {r} \leqslant n} f (e _ {i _ {1}} \otimes \dots \otimes e _ {i _ {r}}) (e _ {i _ {1}} ^ {*} \otimes \dots \otimes e _ {i _ {r}} ^ {*}).
\]

\begin{enumerate}
\def\labelenumi{(\alph{enumi})}
\setcounter{enumi}{1}
\tightlist
\item
  Let \(\beta\) be a non-degenerate bilinear form on \(\mathbf{M}\) which is invariant under \(\mathfrak{g}\) . Let \((e_1', \ldots, e_n')\) be the basis of \(\mathbf{M}\) such that \(\beta(e_i, e_j') = \delta_{ij}\) . Let \(f\) be an invariant linear form on \(\bigotimes^r \mathbf{M}\) . Deduce from (a) that
\end{enumerate}

\[
\sum_ {1 \leqslant i _ {1}, \dots , i _ {r} \leqslant n} f (e _ {i _ {1}} \otimes \dots \otimes e _ {i _ {r}}) (e _ {i _ {1}} ^ {\prime} \otimes \dots \otimes e _ {i _ {r}} ^ {\prime})
\]

is an invariant element of \(\otimes\) M independent of the choice of basis \((e_i)\) .

\begin{enumerate}
\def\labelenumi{(\alph{enumi})}
\setcounter{enumi}{2}
\item
  Let U be the enveloping algebra of g and U* the dual space of U. The adjoint representation of g can be extended to a representation \(x \mapsto x_{\mathrm{U}}\) of g on U defined by \(x_{\mathrm{U}}u = xu - ux\) for all \(u \in U\) . Hence U* has a g-module structure. For an element f of U* to be invariant, it is necessary and sufficient that \(f(uv) = f(vu)\) for all u, v in U.
\item
  Let \(\mathfrak{h}\) be a finite-dimensional ideal of \(\mathfrak{g}\) and \(\gamma\) an invariant bilinear form on \(\mathfrak{g}\) whose restriction to \(\mathfrak{h}\) is non-degenerate. Let \((y_i)_{1\leq i\leq n}, (y_j')_{1\leq j\leq n}\) be two bases of \(\mathfrak{h}\) such that \(\gamma(y_i,y_j') = \delta_{ij}\) . Let \(r\) be an integer \(\geqslant 1\) . Let \(f\) be a linear form on U such that \(f(w) = f(vu)\) for all \(u,v\) in U. Deduce from (b) and (c) that
\end{enumerate}

\[
\sum_ {1 \leqslant t _ {1}, \dots t _ {r} \leqslant n} f (y _ {t _ {1}} y _ {t _ {2}} \cdot \cdot \cdot y _ {t _ {r}}) y _ {t _ {1}} ^ {\prime} y _ {t _ {2}} ^ {\prime} \cdot \cdot \cdot y _ {t _ {r}} ^ {\prime}
\]

is an element of the centre of U independent of the choice of basis \((y_{i})\) . In particular recover the Casimir elements.

\begin{enumerate}
\def\labelenumi{(\alph{enumi})}
\setcounter{enumi}{4}
\tightlist
\item
  Let \(\theta\) be a permutation of \(\{1, \ldots, r\}\) . Arguing similarly, show that
\end{enumerate}

\[
\sum_ {1 \leqslant i _ {1}, \dots i _ {r} \leqslant n} f (y _ {i _ {\theta (1)}} y _ {i _ {\theta (2)}} \dots y _ {i _ {\theta (r)}}) y _ {i _ {1}} ^ {\prime} y _ {i _ {2}} ^ {\prime} \dots y _ {i _ {r}} ^ {\prime}
\]

is an element of the centre of U independent of the choice of basis \((y_{i})\) .

\begin{enumerate}
\def\labelenumi{\arabic{enumi}.}
\setcounter{enumi}{6}
\item
  Let g be a complex Lie algebra, \(\beta\) its Killing form and \(g_{0}\) the Lie algebra derived from g by restricting the base field to R. Show that the Killing form of \(g_{0}\) is twice the real part of \(\beta\) .
\item
  Let \(\mathfrak{g}\) be a real Lie algebra. The multilinear form
\end{enumerate}

\[
(x _ {1}, \dots , x _ {n}) \mapsto \operatorname{Tr} (\operatorname{ad} x _ {1} \circ \operatorname{ad} x _ {2} \circ \dots \circ \operatorname{ad} x _ {n})
\]

on \(\mathfrak{g}^n\) is invariant.

\begin{enumerate}
\def\labelenumi{\arabic{enumi}.}
\setcounter{enumi}{8}
\item
  Let g be a Lie K-algebra, \(\rho\) and \(\rho'\) semi-simple representations of g on K-modules V, \(V'\) and \(\phi\) a homomorphism of the g-module V onto the g-module \(V'\) . Show that the image under \(\phi\) of the set of invariants of V is the set of invariants of \(V'\) (use Proposition 6).
\item
  Suppose that K is a field. Let g be a Lie K-algebra and \(M_{1}\) , \(M_{2}\) two non-isomorphic simple g-modules of finite dimension over K. If \(K_{1}\) is a separable extension of K, show that there exists no simple \(g_{(K_{1})}\) -module isomorphic both to a sub- \(g_{(K_{1})}\) -module of \(M_{1(K_{1})}\) and to a sub- \(g_{(K_{1})}\) -module of \(M_{2(K_{1})}\) . (Use Proposition 4 and note that the existence of an invariant element \(\neq0\) in \(M_{1(K_{1})}^{*}\otimes_{K}M_{2(K_{1})}\) , implies the existence of an invariant element \(\neq0\) in \(M_{1}^{*}\otimes_{K}M_{2}\) (Algebra, Chapter II, § 5, no. 3, Theorem 1.)
\item
  Let \(\mathfrak{w}\) be the Lie \(p\) -algebra defined in Exercise 21 of § 1. Show that the Killing form of \(\mathfrak{w}\) is zero.
\end{enumerate}

¶ 12. Suppose that K is a field. Let g be a Lie K-algebra and M a g-module. Let \(C^{p}(\mathfrak{g}, \mathrm{M})\) denote the space of alternating linear mappings of \(g^{p}\) into M. We write \(C^{0}(\mathfrak{g}, \mathrm{M}) = \mathrm{M}\) and, for p \textless{} 0, \(C^{p}(\mathfrak{g}, \mathrm{M}) = \{0\}\) . Let \(C^{*}(\mathfrak{g}, \mathrm{M})\) be the direct sum of the \(C^{p}(\mathfrak{g}, \mathrm{M})\) . The elements of \(C^{*}(\mathfrak{g}, \mathrm{M})\) are called the cochains of g with values in M; those of \(C^{p}(\mathfrak{g}, \mathrm{M})\) are said to be of degree p.~For all \(y \in g\) , we denote by \(i(y)\) the endomorphism of \(C^{*}(\mathfrak{g}, \mathrm{M})\) which maps each subspace \(C^{p}(\mathfrak{g}, \mathrm{M})\) into \(C^{p-1}(\mathfrak{g}, \mathrm{M})\) and which, for p \textgreater{} 0, is given by the formula:

\[
(i (y) f) \left(x _ {1}, \dots , x _ {p - 1}\right) = f \left(y, x _ {1}, \dots , x _ {p - 1}\right).\tag{1}
\]

We have \(i(y)^{2}=0\) .

\begin{enumerate}
\def\labelenumi{(\alph{enumi})}
\tightlist
\item
  The adjoint representation of g and the representation of g on M define a representation of g on the space of multilinear mappings of \(g^{p}\) into M. Show that \(C^{p}(g, M)\) is stable under this representation. Let \(\theta\) be the representation of g on \(C^{*}(g, M)\) thus defined. Show that:
\end{enumerate}

\[
\theta (x) i (y) - i (y) \theta (x) = i ([ x, y ])
\]

for all x, y in g.

\begin{enumerate}
\def\labelenumi{(\alph{enumi})}
\setcounter{enumi}{1}
\tightlist
\item
  Show that there exists one and only one endomorphism \(d\) of \(\mathrm{C}^*(\mathfrak{g},\mathrm{M})\) , mapping \(\mathrm{C}^p (\mathfrak{g},\mathrm{M})\) into \(\mathrm{C}^{p + 1}(\mathfrak{g},\mathrm{M})\) , such that
\end{enumerate}

\[
d i (y) + i (y) d = \theta (y)\tag{3}
\]

for all \(y \in \mathfrak{g}\) . (Argue by induction on the degree of the cochains.) Show that, for \(f \in \mathrm{C}^p(\mathfrak{g}, \mathbf{M})\) :

\[
\begin{array}{r l} d f (x _ {1}, x _ {2}, \dots , x _ {p + 1}) & = \sum_ {i <   j} (- 1) ^ {i + j} f ([ x _ {i}, x _ {j} ], x _ {1}, \dots , \hat {x} _ {i}, \dots , \hat {x} _ {j}, \dots , x _ {p + 1}) \\ & \quad + \sum_ {i} (- 1) ^ {i + 1} (x _ {i}) _ {\mathrm{M}} f (x _ {1}, \dots , \hat {x} _ {i}, \dots , x _ {p + 1}) \end{array}
\]

(where the symbol \^{} over a letter means that it is omitted).

\begin{enumerate}
\def\labelenumi{(\alph{enumi})}
\setcounter{enumi}{2}
\tightlist
\item
  Show that, for all \(y \in g\) ,
\end{enumerate}

\[
d \theta (y) = \theta (y) d.\tag{4}
\]

(Show first, using (2) and (3), that \(d\theta(y) - \theta(y)d\) commutes with every \(i(x)\) . Then argue by induction on the degree of the cochains.)

\begin{enumerate}
\def\labelenumi{(\alph{enumi})}
\setcounter{enumi}{3}
\item
  Show that \(d^2 = 0\) . (Show first, using (3) and (4), that \(d^2\) commutes with every \(i(x)\) . The argue by induction on the degree of the cochains.)
\item
  The restriction of \(d\) to \(\mathrm{C}^p (\mathfrak{g},\mathrm{M})\) has a kernel \(\mathrm{Z}^p (\mathfrak{g},\mathrm{M})\) whose elements are called cocycles of degree \(p\) with values in M. The restriction of \(d\) to \(\mathrm{C}^{p - 1}(\mathfrak{g},\mathrm{M})\) has image \(\mathrm{B}^p (\mathfrak{g},\mathrm{M})\) , whose elements are called coboundaries of degree \(p\) with values in M. Then \(\mathrm{B}^p (\mathfrak{g},\mathrm{M})\subset \mathrm{Z}^p (\mathfrak{g},\mathrm{M})\) . The quotient space
\end{enumerate}

\[
\mathrm{Z} ^ {p} (\mathfrak {g}, \mathrm{M}) / \mathrm{B} ^ {p} (\mathfrak {g}, \mathrm{M}) = \mathrm{H} ^ {p} (\mathfrak {g}, \mathrm{M})
\]

 is called the cohomology space of g of degree p with values in M. The direct sum of the \(\mathrm{H}^{p}(\mathfrak{g},\mathrm{M})\) is denoted by \(\mathrm{H}^{*}(\mathfrak{g},\mathrm{M})\) . Show that \(\mathrm{H}^{0}(\mathfrak{g},\mathrm{M})\) is identified with the set of invariants of M. Let \(\phi\) be a homomorphism of the g-module M into the g-module N. For all \(f \in \mathrm{C}^{p}(\mathfrak{g},\mathrm{M})\) , \(\phi \circ f \in \mathrm{C}^{p}(\mathfrak{g},\mathrm{N})\) , so that \(\phi\) can be extended to a K-linear mapping \(\phi'\) : \(\mathrm{C}^{*}(\mathfrak{g},\mathrm{M}) \to \mathrm{C}^{*}(\mathfrak{g},\mathrm{N})\) . Show that \(\phi' \circ d = d \circ \phi'\) . Deduce that \(\phi'\) defines a homomorphism

\[
\tilde {\phi} \colon \mathrm{H} ^ {*} (\mathfrak {g}, \mathrm{M}) \to \mathrm{H} ^ {*} (\mathfrak {g}, \mathrm{N})
\]

which is said to be associated with \(\phi\) .

\begin{enumerate}
\def\labelenumi{(\alph{enumi})}
\setcounter{enumi}{5}
\tightlist
\item
  Let L be a sub-g-module of M and N the quotient g-module M/L. The canonical homomorphisms \(L \stackrel{\iota}{\rightarrow} M \stackrel{p}{\rightarrow} N\) define homomorphisms
\end{enumerate}

\[
\begin{array}{l} \mathrm{C} ^ {*} (\mathfrak {g}, \mathrm{L}) \xrightarrow {\iota^ {\prime}} \mathrm{C} ^ {*} (\mathfrak {g}, \mathrm{M}) \xrightarrow {p ^ {\prime}} \mathrm{C} ^ {*} (\mathfrak {g}, \mathrm{N}), \\ \mathrm{H} ^ {n} (\mathfrak {g}, \mathrm{L}) \xrightarrow {\iota^ {n}} \mathrm{H} ^ {n} (\mathfrak {g}, \mathrm{M}) \xrightarrow {p ^ {n}} \mathrm{H} ^ {n} (\mathfrak {g}, \mathrm{N}). \end{array}
\]

Show that \(\iota'\) is injective and has image the kernel of \(p'\) and that \(p'\) is surjective. For all \(c \in \mathrm{H}^n(\mathfrak{g},\mathrm{N})\) , let \(z \in \mathrm{Z}^n(\mathfrak{g},\mathrm{N})\) be a representative of \(c\) and \(a \in \mathrm{C}^n(\mathfrak{g},\mathrm{M})\) be such that \(p'(a) = z\) ; show that \(da \in \mathrm{Z}^{n+1}(\mathfrak{g},\mathrm{L})\) and that its class in \(\mathrm{H}^{n+1}(\mathfrak{g},\mathrm{L})\) depends only on \(c\) ; denote this class by \(\delta^n c\) . Show that the sequence of homomorphisms:

\[
\begin{array}{r l} 0 \longrightarrow & \mathrm{H} ^ {0} (\mathfrak {g}, \mathrm{L}) \stackrel {{\iota ^ {0}}} {{\longrightarrow}} \mathrm{H} ^ {0} (\mathfrak {g}, \mathrm{M}) \stackrel {{p ^ {0}}} {{\longrightarrow}} \mathrm{H} ^ {0} (\mathfrak {g}, \mathrm{N}) \stackrel {{\delta^ {0}}} {{\longrightarrow}} \\ & \longrightarrow \mathrm{H} ^ {1} (\mathfrak {g}, \mathrm{L}) \stackrel {{\iota ^ {1}}} {{\longrightarrow}} \mathrm{H} ^ {1} (\mathfrak {g}, \mathrm{M}) \stackrel {{p ^ {1}}} {{\longrightarrow}} \mathrm{H} ^ {1} (\mathfrak {g}, \mathrm{N}) \stackrel {{\delta^ {1}}} {{\longrightarrow}} \dots \end{array}
\]

is an exact sequence.

\begin{enumerate}
\def\labelenumi{(\alph{enumi})}
\setcounter{enumi}{6}
\tightlist
\item
  With the notation of (f), the exact sequence:
\end{enumerate}

\[
0 \to \mathscr {L} (\mathrm{N}, \mathrm{L}) \to \mathscr {L} (\mathrm{N}, \mathrm{M}) \to \mathscr {L} (\mathrm{N}, \mathrm{N}) \to 0
\]

defines an exact sequence

\[
\mathrm{H} ^ {0} (\mathfrak {g}, \mathscr {L} (\mathrm{N}, \mathrm{M})) \rightarrow \mathrm{H} ^ {0} (\mathfrak {g}, \mathscr {L} (\mathrm{N}, \mathrm{N})) \rightarrow \mathrm{H} ^ {1} (\mathfrak {g}, \mathscr {L} (\mathrm{N}, \mathrm{L})).
\]

 The identity mapping of N is an invariant u of \(\mathcal{L}(\mathrm{N},\mathrm{N})\) and hence an element of \(\mathrm{H}^{0}(\mathfrak{g},\mathcal{L}(\mathrm{N},\mathrm{N}))\) ; let c be its image in \(\mathrm{H}^{1}(\mathfrak{g},\mathcal{L}(\mathrm{N},\mathrm{L}))\) . Then for the existence in M of a sub-g-module supplementary to L it is necessary and sufficient that c=0. (This condition means that u is the image of an element of \(\mathrm{H}^{0}(\mathfrak{g},\mathcal{L}(\mathrm{N},\mathrm{M}))\) , that is a homomorphism v of the g-module N into the g-module M; then \(v(\mathrm{N})\) is the desired supplement.)

\begin{enumerate}
\def\labelenumi{(\alph{enumi})}
\setcounter{enumi}{7}
\item
  Show that \(Z^{1}(\mathfrak{g},\mathfrak{g})\) (where g is considered as a g-module by means of the adjoint representation) is identified with the vector space of derivations of g and that \(B^{1}(\mathfrak{g},\mathfrak{g})\) is identified with the vector space of inner derivations of g.
\item
  Let \(\mathfrak{a}\) and \(\mathfrak{b}\) be two Lie K-algebras with \(\mathfrak{b}\) commutative and \(\mathfrak{b} \xrightarrow{\lambda} \mathfrak{g} \xrightarrow{\mu} \mathfrak{a}\) an extension of \(\mathfrak{a}\) by \(\mathfrak{b}\) . For all \(x \in \mathfrak{g}\) , the restriction of \(\mathrm{ad}_{\mathfrak{a}} x\) to \(\mathfrak{b}\) depends only on the class of \(x\) modulo \(\mathfrak{b}\) , whence there is an \(\mathfrak{a}\) -module structure on \(\mathfrak{b}\) . Let \(\nu\) be a K-linear mapping of \(\mathfrak{a}\) into \(\mathfrak{g}\) such that \(\mu \circ \nu\) is the identity mapping of \(\mathfrak{a}\) . For \(x, y\) in \(\mathfrak{a}\) , we write \(f(x, y) = [\nu x, \nu y] - \nu([x, y])\) . Show that \(f \in Z^2(\mathfrak{a}, \mathfrak{b})\) and that the class \(c\) of \(f\) in \(H^2(\mathfrak{a}, \mathfrak{b})\) does not depend on the choice of \(\nu\) . For two extensions of \(\mathfrak{a}\) by \(\mathfrak{b}\) , which define the same \(\mathfrak{a}\) -module structure on \(\mathfrak{b}\) , to be equivalent, it is necessary and sufficient that the corresponding class \(c \in H^2(\mathfrak{a}, \mathfrak{b})\) be the same. For the extension to be inessential, it is necessary and sufficient that \(c = 0\) . If B is an \(\mathfrak{a}\) -module and \(c \in H^2(\mathfrak{a}, \mathfrak{B})\) , there exists an extension of \(\mathfrak{a}\) by B (considered as a commutative Lie algebra) defining the given \(\mathfrak{a}\) -module structure on B and the given element \(c\) of \(H^2(\mathfrak{a}, \mathfrak{B})\) .
\item
  Let g be a finite-dimensional Lie algebra over K, U its enveloping algebra, h an ideal of g, β an invariant bilinear form on g whose restriction to h is non-degenerate, \((y_{i})_{1 \leqslant i \leqslant n}\) and \((y_{i}')_{1 \leqslant i \leqslant n}\) two bases of h such that
\end{enumerate}

 \(\beta(y_i, y_i') = \delta_{ij}, t = \sum_{i=1}^{n} y_i y_i' \in \mathbf{U}, \mathbf{M}\) a \(g\) -module and \(\rho\) the endomorphism of \(\mathrm{C}^*(\mathfrak{g}, \mathbf{M})\) which maps each \(\mathrm{C}^p(\mathfrak{g}, \mathbf{M})\) into \(\mathrm{C}^{p-1}(\mathfrak{g}, \mathbf{M})\) and which, for \(p > 0\) , is defined by:

\[
(\wp f) (x _ {1}, \dots , x _ {p - 1}) = \sum_ {k = 1} ^ {n} (y _ {k}) _ {\mathrm{M}} f (y _ {k} ^ {\prime}, x _ {1}, \dots , x _ {p - 1}).
\]

Finally let \(\Gamma\) be the endomorphism of \(C^*(g, M)\) extending \(t_M\) which maps each cochain \(f\) of degree \(>0\) onto the cochain \(t_M \circ f\) . Show that \(\rho d + d\rho = \Gamma\) . (If for \(x \in g\) we write:

\[
[ x, y _ {k} ] = \sum_ {l = 1} ^ {n} a _ {k l} y _ {l}, \quad [ x, y _ {k} ^ {\prime} ] = \sum_ {l = 1} ^ {n} a _ {k l} ^ {\prime} y _ {l} ^ {\prime},
\]

·show first that \(a_{kl} = -a_{lk}^{\prime}\) .) Deduce that \(\Gamma d = d\Gamma\) . If M is simple and finite-dimensional over K, \(\beta\) is the associated bilinear form and dim \(\mathfrak{h}\) is not divisible by the characteristic of K, show that \(\mathrm{H}^{p}(\mathfrak{g},\mathrm{M})=\{0\}\) for all p.~(Using Proposition 12, show that \(\Gamma\) is an automorphism of \(\mathrm{C}^{*}(\mathfrak{g},\mathrm{M})\) and hence induces automorphisms of \(\mathrm{Z}^{p}(\mathfrak{g},\mathrm{M})\) and \(\mathrm{B}^{p}(\mathfrak{g},\mathrm{M})\) . For

\[
f \in \mathrm{Z} ^ {p} (\mathfrak {g}, \mathrm{M}), \quad \Gamma f = (d \rho + \rho d) f = d \rho f \in \mathrm{B} ^ {p} (\mathfrak {g}, \mathrm{M}),
\]

whence \(Z^{p}(\mathfrak{g},\mathbf{M}) = \mathrm{B}^{p}(\mathfrak{g},\mathbf{M}).\)

\exercisesfor{4}

The conventions of § 4 remain valid unless otherwise mentioned.

\begin{enumerate}
\def\labelenumi{\arabic{enumi}.}
\item
  Let \(\mathfrak{g}\) be a nilpotent Lie algebra and \(p\) (resp. \(q\) ) the smallest integer such that \(\mathcal{C}^p\mathfrak{g} = \{0\}\) (resp. \(\mathcal{C}_q\mathfrak{g} = \mathfrak{g}\) ). Show that \(p = q + 1\) and that \(\mathcal{C}_i\mathfrak{g} \supset \mathcal{C}^{p - i}\mathfrak{g}\) . (Use the argument of Proposition 1.)
\item
  Let \(\mathfrak{g}\) be a semi-direct product of an algebra \(\mathfrak{h}\) of dimension 1 and a commutative ideal \(\mathfrak{g}'\) . Let \(x \in \mathfrak{h}, x \neq 0\) , and \(u\) be the restriction of \(\operatorname{ad}_{\mathfrak{g}} x\) to \(\mathfrak{g}'\) .
\end{enumerate}

\begin{enumerate}
\def\labelenumi{(\alph{enumi})}
\item
  For g to be nilpotent, it is necessary and sufficient that u be nilpotent.
\item
  For the Killing form of \(\mathfrak{g}\) to be zero, it is necessary and sufficient that \(\operatorname{Tr}(u^2) = 0\) .
\item
  Deduce from (a) and (b) that there exist non-nilpotent Lie algebras whose Killing forms are zero.
\item
  Deduce from (a) that in a nilpotent Lie algebra such that \(\mathcal{C}^{p-1}\mathfrak{g} \neq \{0\}\) , \(\mathcal{C}^p\mathfrak{g} = \{0\}\) it is possible that \(\mathcal{C}_i\mathfrak{g} \neq \mathcal{C}^{p-i}\mathfrak{g}\) .
\end{enumerate}

\begin{enumerate}
\def\labelenumi{\arabic{enumi}.}
\setcounter{enumi}{2}
\tightlist
\item
  \begin{enumerate}
  \def\labelenumii{(\alph{enumii})}
  \tightlist
  \item
    Let \(\mathfrak{g}\) be a nilpotent Lie algebra, \(\mathfrak{z}\) its centre and \(\mathfrak{b}\) a non-zero ideal of \(\mathfrak{g}\) . Show that \(\mathfrak{z} \cap \mathfrak{b} \neq \{0\}\) . (Consider \(\mathfrak{b}\) as a \(\mathfrak{g}\) -module by means of the adjoint representation.)
  \end{enumerate}
\end{enumerate}

\begin{enumerate}
\def\labelenumi{(\alph{enumi})}
\setcounter{enumi}{1}
\tightlist
\item
  If in a Lie algebra \(\mathfrak{g}\) an ideal \(\mathfrak{b}\) is contained in \(\mathcal{C}_{i+1}\mathfrak{g}\) but not in \(\mathcal{C}_i\mathfrak{g}\) , show that the ideals \(\mathfrak{b} \cap \mathcal{C}_k\mathfrak{g}\) are all distinct for \(0 \leqslant k \leqslant i + 1\) . (Apply (a).)
\end{enumerate}

\begin{enumerate}
\def\labelenumi{\arabic{enumi}.}
\setcounter{enumi}{3}
\tightlist
\item
  Let g be a nilpotent Lie algebra.
\end{enumerate}

\begin{enumerate}
\def\labelenumi{(\alph{enumi})}
\item
  Every subalgebra of \(\mathfrak{g}\) is subinvariant. (Use Proposition 3.)
\item
  Let \(\mathfrak{b}\) be a vector subspace of \(\mathfrak{g}\) such that \(\mathfrak{b} + \mathcal{D}\mathfrak{g} = \mathfrak{g}\) . Show that the subalgebra of \(\mathfrak{g}\) generated by \(\mathfrak{b}\) is \(\mathfrak{g}\) . (Apply (a) to this subalgebra. Reduce it thus to the case where \(\mathfrak{b}\) is an ideal of \(\mathfrak{g}\) and use Proposition 4 of § 1.) Deduce that the minimum number of generators of \(\mathfrak{g}\) is \(\dim \mathfrak{g}/\mathcal{D}\mathfrak{g}\) .
\end{enumerate}

\begin{enumerate}
\def\labelenumi{\arabic{enumi}.}
\setcounter{enumi}{4}
\tightlist
\item
  \begin{enumerate}
  \def\labelenumii{(\alph{enumii})}
  \tightlist
  \item
    Show that a Lie algebra g in which every subalgebra is subinvariant is nilpotent. (Show that if \(\dim g > 1\) every \(x \in g\) belongs to an ideal \(h \neq g\) of g which has the same property as g and hence is nilpotent by induction on the dimension of g; hence \(ad_{g} x\) is nilpotent.)
  \end{enumerate}
\end{enumerate}

\begin{enumerate}
\def\labelenumi{(\alph{enumi})}
\setcounter{enumi}{1}
\tightlist
\item
  Show that if in a Lie algebra \(\mathfrak{g}\) every subalgebra distinct from \(\mathfrak{g}\) is distinct from its normalizer, \(\mathfrak{g}\) is nilpotent. (Reduce it to (a).)
\end{enumerate}

\begin{enumerate}
\def\labelenumi{\arabic{enumi}.}
\setcounter{enumi}{5}
\item
  Let g be a nilpotent Lie algebra and a a commutative ideal of g. The following conditions are equivalent: (a) a is a maximal commutative ideal of g; (b) a is a maximal commutative subalgebra of g: (c) a is equal to its centralizer \(a'\) in g. (The implications not (a) \(\Rightarrow\) not (b) \(\Rightarrow\) not (c) are obvious. If \(a' \neq a\) , there exists by Proposition 1 an ideal \(a''\) of g such that \(a \subset a'' \subset a'\) and dim \(a''/a = 1\) . Then \(a'' = a + Kx\) , hence \([a'', a''] \subset [a, a] + [x, a] = \{0\}\) and hence (a) is false.)
\item
  \begin{enumerate}
  \def\labelenumii{(\alph{enumii})}
  \tightlist
  \item
    Let V be a finite-dimensional vector space over \(K\), \(\mathfrak{g}\) a Lie algebra of nilpotent endomorphisms of V and \((\mathrm{V}_{r})_{0 \leq r \leq n}\) a decreasing sequence of vector subspaces of V with \(V_{0} = V, V_{n} = \{0\}\) and \(\mathfrak{g}(\mathrm{V}_{r}) \subset V_{r+1}\) for \(0 \leq r < n\) . Show by induction on i that \((\mathcal{D}^{i}\mathfrak{g})(\mathrm{V}_{r}) \subset V_{r+2^{i}}\) . If \(\dim V \geq 2^{i}\) , the subspace of elements of V annihilated by \(\mathcal{D}^{i}\mathfrak{g}\) is of dimension \(\geqslant 2^{i}\) . If \(\dim V \leqslant 2^{i}\) , \(\mathcal{D}^{i}\mathfrak{g} = \{0\}\) .
  \end{enumerate}
\end{enumerate}

\begin{enumerate}
\def\labelenumi{(\alph{enumi})}
\setcounter{enumi}{1}
\item
  Let \(\mathfrak{g}\) be a nilpotent Lie algebra and \(i\) an integer \(\geqslant 0\) . If \(\dim \mathcal{D}^{i}\mathfrak{g} > 2^{i} + 1\) the centre of \(\mathcal{D}^{i}\mathfrak{g}\) is of dimension \(\geqslant 2^{i}\) . If \(\dim \mathcal{D}^{i}\mathfrak{g} \leqslant 2^{i} + 1\) , \(\mathcal{D}^{i}\mathfrak{g}\) is commutative. (Apply (a) to the restrictions to \(\mathcal{D}^{i}\mathfrak{g}\) of the \(\operatorname{ad}_{\mathfrak{g}} x, x \in \mathfrak{g}\) .)
\item
  Let \(\mathfrak{g}\) be a nilpotent Lie algebra and \(i\) an integer \(\geqslant 0\) . If \(\mathcal{D}^{i}\mathfrak{g}\) is not commutative, \(\mathcal{D}^{i}\mathfrak{g}/\mathcal{D}^{i+1}\mathfrak{g}\) is of dimension \(\geqslant 2^{i} + 1\) . (Reduce it, by passing to the quotient, to the case where \(\dim \mathcal{D}^{i+1}\mathfrak{g} = 1\) and use (b).)
\end{enumerate}

\begin{enumerate}
\def\labelenumi{\arabic{enumi}.}
\setcounter{enumi}{7}
\tightlist
\item
  Let \(\mathfrak{g}\) be a Lie algebra, \(\mathfrak{m}\) an ideal of codimension 1, \(x\) an element of \(\mathfrak{g}\) not belonging to \(\mathfrak{m}\) and \(z \in \mathfrak{g}\) .
\end{enumerate}

\begin{enumerate}
\def\labelenumi{(\alph{enumi})}
\item
  Show that the linear mapping D which reduces to 0 on m and maps x to z is a derivation if z belongs to the centralizer a of m in g.
\item
  Let \(q\) be the largest integer such that \(a \subset \mathcal{C}^q\mathfrak{g}\) . Show that if further \(z \notin \mathcal{C}^{q+1}\mathfrak{g}\) , \(D\) is not an inner derivation of \(\mathfrak{g}\) .
\item
  Deduce from (a) and (b) and Exercise 7 (c) that if g is a nilpotent Lie algebra of dimension \textgreater1, the vector space of inner derivations of g is of co-dimension \(\geqslant2\) in the space of all derivations of g.
\end{enumerate}

\begin{enumerate}
\def\labelenumi{\arabic{enumi}.}
\setcounter{enumi}{8}
\tightlist
\item
  \begin{enumerate}
  \def\labelenumii{(\alph{enumii})}
  \tightlist
  \item
    Verify that the following multiplication tables define two nilpotent Lie algebras \(\mathfrak{g}_3, \mathfrak{g}_4\) of dimensions 3 and 4:
  \end{enumerate}
\end{enumerate}

\[
\mathfrak {g} _ {3}: \left[ x _ {1}, x _ {2} \right] = x _ {3}, \quad \left[ x _ {1}, x _ {3} \right] = \left[ x _ {2}, x _ {3} \right] = 0
\]

\[
\mathfrak {g} _ {4}: \left[ x _ {1}, x _ {2} \right] = x _ {3}, \quad \left[ x _ {1}, x _ {3} \right] = x _ {4}, \quad \left[ x _ {1}, x _ {4} \right] = \left[ x _ {2}, x _ {3} \right] = \left[ x _ {2}, x _ {4} \right] = \left[ x _ {3}, x _ {4} \right] = 0.
\]

\begin{enumerate}
\def\labelenumi{(\alph{enumi})}
\setcounter{enumi}{1}
\item
  Show that \(\mathfrak{g}_3\) is isomorphic to \(\mathfrak{n}(3, \mathbf{K})\) and also to \(\mathfrak{sl}(2, \mathbf{K})\) if \(\mathbf{K}\) is of characteristic 2.
\item
  Let \(\mathfrak{g}_1\) be the unique Lie algebra of dimension 1. Show that the nilpotent Lie algebras of dimension \(\leqslant 4\) are given by the following table:
\end{enumerate}

\[
\text { dimension } 1: \mathfrak {g} _ {1}; \text { dimension } 2: (\mathfrak {g} _ {1}) ^ {2};
\]

dimension 3: \((\mathfrak{g}_1)^3, \mathfrak{g}_3\) ; dimension 4: \((\mathfrak{g}_1)^4, \mathfrak{g}_3 \times \mathfrak{g}_1, \mathfrak{g}_4\) .

(Use Exercise 7 (c). If \(\mathfrak{g}\) is nilpotent and 3-dimensional and \(\dim \mathcal{D}\mathfrak{g} = 1\) , observe that \(\mathcal{D}\mathfrak{g}\) is contained in the centre of \(\mathfrak{g}\) (Exercise 3 (a)), whence \(\mathfrak{g} = \mathfrak{g}_3\) . If dim \(\mathfrak{g} = 4\) , dim \(\mathcal{D}\mathfrak{g} = 1\) , observe that the bracket on \(\mathfrak{g}\) defines an alternating bilinear form on \(\mathfrak{g} / \mathcal{D}\mathfrak{g}\) , which is necessarily degenerate and hence that there exists a subalgebra \(\mathfrak{h}\) of \(\mathfrak{g}\) contained in the centre of \(\mathfrak{g}\) with dim \(\mathfrak{h} = 1\) , \(\mathfrak{h}\cap \mathcal{D}\mathfrak{g} = \{0\}\) , whence \(\mathfrak{g} = \mathfrak{h}\times \mathfrak{g}_3\) . If dim \(\mathfrak{g} = 4\) , dim \(\mathcal{D}\mathfrak{g} = 2\) , \(\mathcal{D}\mathfrak{g}\) is commutative; applying Theorem 1 to the restrictions to \(\mathcal{D}\mathfrak{g}\) of the \(\operatorname{ad}_{\mathfrak{g}} x\) \((x\in \mathfrak{g})\) , show that there exists a commutative ideal \(\mathfrak{h}\) of \(\mathfrak{g}\) with \(\mathfrak{h}\supseteq \mathcal{D}\mathfrak{g}\) , dim \(\mathfrak{h} = 3\) ; let \(x\in \mathfrak{g},\) \(x\notin \mathfrak{h}\) ; choose a basis of \(\mathfrak{h}\) such that the restriction of \(\operatorname{ad}_{\mathfrak{g}}x\) to \(\mathfrak{h}\) has a Jordan matrix with respect to this basis; then \(\mathfrak{g} = \mathfrak{g}_4\).

\begin{enumerate}
\def\labelenumi{\arabic{enumi}.}
\setcounter{enumi}{9}
\tightlist
\item
  Let \(\mathcal{D}\) be the Lie algebra of derivations of the algebra \(\mathfrak{g}_4\) (Exercise 9). Show that \(\mathcal{D}\) is of dimension 7 and of zero centre and that the ideal \(\mathcal{D}' \subset \mathcal{D}\) of inner derivations has no supplementary subalgebra in \(\mathcal{D}\) .
\end{enumerate}

¶ 11. Let L be a commutative ring, A a left Artinian algebra over L and γ a mapping of A × A into L. Define on A the internal law of composition \((a, b) \mapsto a * b = ab + \gamma(a, b)ba\) . For every subset E of A, let \(\tilde{E}\) denote the subring (with no unit element) of A generated by E. Show that if E consists of nilpotent elements and is stable under the law *, then \(\tilde{E}\) is nilpotent (that is, there exists n \textgreater{} 0 such that \(\tilde{E}^{n} = \{0\}\) ). Proceed as follows:

\begin{enumerate}
\def\labelenumi{\arabic{enumi}.}
\item
  Suppose first that A is a simple ring and hence isomorphic to \(\mathcal{L}_{\mathrm{D}}(\mathrm{T})\) , where T is a left vector space of finite dimension \(m\) over a field D. Argue by induction on \(m\) . Let \(\Phi\) be the set of subsets \(\mathrm{F} \subset \mathrm{E}\) which are stable under \(*\) and such that \(\tilde{\mathrm{F}}\) is nilpotent. Show that \(\tilde{\Phi}\) admits a maximal element M (note that \(\tilde{\mathrm{F}}^m = \{0\}\) for all \(\mathrm{F} \in \Phi\) ). Suppose \(\tilde{\mathbf{M}} \neq \tilde{\mathbf{E}}\) . Show that there exists \(a \in \mathrm{E}\) such that \(a \notin \tilde{\mathbf{M}}\) and \(t * a \in \tilde{\mathbf{M}}\) for all \(t \in \mathbf{M}\) (observe that there can be no infinite sequence \((a_n)\) such that \(a_n \in \mathrm{E}\) , \(a_n \notin \mathbf{M}\) , \(a_n = t_{n-1} * a_{n-1}\) with \(t_{n-1} \in \mathbf{M}\) ). Let S be the subspace of T the direct sum of the \(u(\mathrm{T})\) with \(u \in \tilde{\mathbf{M}}\) ; show that \(S \neq \{0\}\) and \(S \neq T\) and that \(a(S) \subset S\) . Let N be the set of \(v \in E\) such that \(v(S) \subset S\) . By using the induction hypothesis and considering the elements of N as operating on S and on T/S, show that \(N \in \Phi\) , which implies a contradiction.
\item
  In the general case, use the fact that the Jacobson radical of A is nilpotent. Deduce from this result a new proof of Engel's Theorem.
\item
  Let \(\mathfrak{g}\) be a Lie algebra and \(\mathcal{C}^{\infty}\mathfrak{g}\) the intersection of the \(\mathcal{C}^p\mathfrak{g}\) .
\end{enumerate}

\begin{enumerate}
\def\labelenumi{(\alph{enumi})}
\item
  The Lie algebra \(\mathfrak{g} / \mathcal{C}^{\infty}\mathfrak{g}\) is nilpotent.
\item
  Show that there exists a nilpotent subalgebra \(\mathfrak{h}\) of \(\mathfrak{g}\) such that \(\mathfrak{g} = \mathfrak{h} + \mathcal{C}^{\infty}\mathfrak{g}\) . (Argue by induction on \(\dim \mathfrak{g}\) . Suppose that \(\mathfrak{g}\) is not nilpotent. Let \(x\in \mathfrak{g}\) be such that \(\mathbf{I} = \bigcap_{k}\left(\mathrm{(ad}x)^{k}(\mathfrak{g})\neq \{0\}\right)\) . Let \(\mathfrak{n}\) be the union of the kernels of the \((\mathrm{ad}x)^k\) . Show that this is a subalgebra of \(\mathfrak{g}\) and that \(\mathfrak{g}\) is the direct sum of \(\mathbf{I}\) and \(\mathfrak{n}\) . By the induction hypothesis \(\mathfrak{n}\) is the sum of \(\mathcal{C}^{\infty}\mathfrak{n}\) and a nilpotent subalgebra \(\mathfrak{h}\) . Finally, \(\mathcal{C}^{\infty}\mathfrak{n}\subset \mathcal{C}^{\infty}\mathfrak{g}\) and \(\mathbf{I}\subset \mathcal{C}^{\infty}\mathfrak{g}\) .
\item
  Verify that the following multiplication table defines a (solvable) Lie algebra g of dimension 5:
\end{enumerate}

\[
\left[ x _ {1}, x _ {2} \right] = x _ {5}, \quad \left[ x _ {1}, x _ {3} \right] = x _ {3}, \quad \left[ x _ {1}, x _ {4} \right] = - x _ {4}, \quad \left[ x _ {3}, x _ {4} \right] = x _ {5}
\]

and the other \([x_i, x_j]\) are zero. Show that for this Lie algebra

\[
\mathcal {C} ^ {\infty} \mathfrak {g} = \mathrm{K} x _ {3} + \mathrm{K} x _ {4} + \mathrm{K} x _ {5}
\]

but that there exists no supplementary subalgebra of \(\mathcal{C}^{\infty}\mathfrak{g}\) in \(\mathfrak{g}\) .

¶ 13. Let g be a Lie algebra and \(\mathfrak{z}\) the centralizer of \(C^{\infty}g\) in g. If \(\mathfrak{z} \not\subset C^{\infty}g\) , g has non-zero centre. (Write \(g = C^{\infty}g + t\) , where t is a nilpotent subalgebra of g (Exercise 12)). Let \(g_{1} = \mathfrak{z} + t\) . Let b be a nilpotent subalgebra such that \(g_{1} = C^{\infty}g_{1} + \mathfrak{b}\) . Show that \(g = C^{\infty}g + \mathfrak{b}\) and that \(C^{\infty}g_{1} \subset \mathfrak{z} \cap C^{\infty}g\) . Let y be an element of \(\mathfrak{z}\) not belonging to \(C^{\infty}g\) . Write \(y = x + x'\) where \(x \in \mathfrak{b}, x' \in C^{\infty}g_{1}\) ; then \(x \in \mathfrak{z}\) and \(x \neq 0\) and hence \(\mathfrak{z} \cap \mathfrak{b}\) is a non-zero ideal of b; using Exercise 3 (a), deduce that there exists a non-zero element of g permutable with \(C^{\infty}g\) and b.)

¶ 14. Let g be a Lie algebra and b a subinvariant subalgebra of g such that the centralizer \(\mathfrak{z}(b)\) of b in g is zero.

\begin{enumerate}
\def\labelenumi{(\alph{enumi})}
\item
  The centralizer \(\mathfrak{z}(\mathcal{C}^{\infty}\mathfrak{b})\) of \(\mathcal{C}^{\infty}\mathfrak{b}\) in \(\mathfrak{g}\) is contained in \(\mathcal{C}^{\infty}\mathfrak{b}\) . (If \(\mathfrak{z}(\mathcal{C}^{\infty}\mathfrak{b})\not\subset\mathcal{C}^{\infty}\mathfrak{b},\mathfrak{z}(\mathcal{C}^{\infty}\mathfrak{b})\not\subset\mathfrak{b}\) by Exercise 13; \(\mathcal{C}^{\infty}\mathfrak{b}\) is an ideal of \(\mathfrak{g}\) (§ 1, Exercise 14); let \(\mathfrak{c}=\mathfrak{b}+\mathfrak{z}(\mathcal{C}^{\infty}\mathfrak{b})\) ; \(\mathfrak{c}\) is a subalgebra, \(\mathfrak{b}\) is subinvariant in \(\mathfrak{c}\) ; \(\mathfrak{b}\) is an ideal of \(\mathfrak{b}_{1}\subset\mathfrak{c}\) , where \(\mathfrak{b}\neq\mathfrak{b}_{1}\) ; let \(y\in\mathfrak{b}_{1},y\notin\mathfrak{b};y=x+x'\) , where \(x\in\mathfrak{z}(\mathcal{C}^{\infty}\mathfrak{b}),x'\in\mathfrak{b}\) ; then \(x\in\mathfrak{z}(\mathcal{C}^{\infty}\mathfrak{b})\cap\mathfrak{b}_{1},x\notin\mathfrak{b};\mathfrak{b}_{1}=\mathrm{Kx}+\mathfrak{b}\) is a subalgebra; \(\mathcal{C}^{\infty}\mathfrak{c}=\mathcal{C}^{\infty}\mathfrak{b}\) because \(\mathfrak{c}\cap\mathfrak{z}(\mathcal{C}^{\infty}\mathfrak{b})\subset\mathcal{C}^{\infty}\mathfrak{b}\) , whence \(\mathcal{C}^{k}\mathfrak{c}\subset\mathcal{C}^{k}\mathfrak{b}\) for all \(k\) ; \(\mathfrak{z}(\mathcal{C}^{\infty}\mathfrak{b})\cap\mathfrak{b}\not\subset\mathcal{C}^{\infty}\mathfrak{b}\) and hence the centre of \(\mathfrak{b}\) is non-zero by Exercise 13; a fortiori, \(\mathfrak{z}(\mathfrak{b})\neq\{0\}\) , which is absurd.)
\item
  Deduce from (a) that, if \(\mathfrak{D}\) denotes the Lie algebra of derivations of \(\mathcal{C}^{\infty}\mathfrak{b}\) and \(\mathfrak{a}\) the centre of \(\mathcal{C}^{\infty}\mathfrak{b}\) , then \(\dim \mathfrak{g} \leqslant \dim \mathfrak{D} + \dim \mathfrak{a}\) (let \(\mathfrak{g}\) operate on \(\mathcal{C}^{\infty}\mathfrak{b}\) by the adjoint representation).
\end{enumerate}

\begin{enumerate}
\def\labelenumi{\arabic{enumi}.}
\setcounter{enumi}{14}
\tightlist
\item
  Let \(l\) be a Lie algebra with zero centre, \(\mathfrak{D}_1\) the Lie algebra of derivations of \(l\) , \(\mathfrak{D}_2\) the Lie algebra of derivations of \(\mathfrak{D}_1, \ldots\) . Then \(l\) is an ideal of \(\mathfrak{D}_1, \mathfrak{D}_1\) is an ideal of \(\mathfrak{D}_2\) , etc. (§ 1, Exercise 15 (c)). Let \(\mathfrak{D}\) be the Lie algebra of derivations of \(\mathcal{C}^\infty l\) and \(a\) the centre of \(\mathcal{C}^\infty l\) .
\end{enumerate}

\begin{enumerate}
\def\labelenumi{(\alph{enumi})}
\item
  Then \(\dim \mathfrak{D}_1 \leqslant \dim \mathfrak{D} + \dim \mathfrak{a}\) . (Use Exercise 14 above and Exercise 15 of § 1.)
\item
  Deduce from (a) that for i sufficiently large all the derivations of \(D_{i}\) are inner.
\end{enumerate}

\begin{enumerate}
\def\labelenumi{\arabic{enumi}.}
\setcounter{enumi}{15}
\item
  Let \(\mathfrak{g}_1\) and \(\mathfrak{g}_2\) be Lie K-algebras and \(\mathfrak{n}_1\) and \(\mathfrak{n}_2\) their largest nilpotent ideals. Show that the largest nilpotent ideal of \(\mathfrak{g}_1 \times \mathfrak{g}_2\) is \(\mathfrak{n}_1 \times \mathfrak{n}_2\) .
\item
  We reconsider the Lie algebra \(\mathfrak{g}_3\) of Exercise 9 (a) with its basis \((x_1, x_2, x_3)\) .
\end{enumerate}

\begin{enumerate}
\def\labelenumi{(\alph{enumi})}
\tightlist
\item
  Let V be the vector space K{[}X{]}. Let D be the differential operator with respect to X on V and M the operator of multiplication by X on V. Show that if K is of characteristic 0 the mapping
\end{enumerate}

\[
\alpha x _ {1} + \beta x _ {2} + \gamma x _ {3} \mapsto \alpha D + \beta M + \gamma
\]

is an infinite-dimensional simple representation \(\rho\) of \(\mathfrak{g}\) on V.

\begin{enumerate}
\def\labelenumi{(\alph{enumi})}
\setcounter{enumi}{1}
\tightlist
\item
  If K is of characteristic p \textgreater{} 0, the ideal \((\mathbf{X}^{p})\) of K{[}X{]} is stable under \(\rho(\mathfrak{g})\) . Under the quotient representation of g on \(\mathrm{K}[\mathrm{X}] / (\mathrm{X}^{p})\) , no line is stable.
\end{enumerate}

\begin{enumerate}
\def\labelenumi{\arabic{enumi}.}
\setcounter{enumi}{17}
\tightlist
\item
  Let \(\mathfrak{g}\) be a 7-dimensional vector space over \(\mathbf{K}\) with a basis \((e_i)_{1\leqslant i\leqslant 7}\) . An alternating bracket is defined on \(\mathfrak{g}\) by the formulae:
\end{enumerate}

\[
\left[ e _ {i}, e _ {j} \right] = \alpha_ {i j} e _ {i + j} \quad (1 \leqslant i <   j \leqslant 7, i + j \leqslant 7)\tag{1}
\]

and all the other brackets \([e_i, e_j]\) are zero for \(i < j\) .

\begin{enumerate}
\def\labelenumi{(\alph{enumi})}
\tightlist
\item
  For the Jacobi identity to hold, it is necessary and sufficient that:
\end{enumerate}

\begin{enumerate}
\def\labelenumi{(\arabic{enumi})}
\setcounter{enumi}{1}
\tightlist
\item
\end{enumerate}

\[
- \alpha_ {2 3} \alpha_ {1 5} + \alpha_ {1 3} \alpha_ {2 4} = 0\tag{2}
\]

\[
\alpha_ {1 2} \alpha_ {3 4} - \alpha_ {2 4} \alpha_ {1 6} + \alpha_ {1 4} \alpha_ {2 5} = 0.\tag{3}
\]

\begin{enumerate}
\def\labelenumi{(\alph{enumi})}
\setcounter{enumi}{1}
\item
  Suppose henceforth that all the \(\alpha_{ij}\) are \(\neq 0\) . Show that the ideals \(\mathcal{C}^2\mathfrak{g}\) , \(\mathcal{C}^3\mathfrak{g}\) , \(\mathcal{C}^4\mathfrak{g}\) , \(\mathcal{C}^5\mathfrak{g}\) , \(\mathcal{C}^6\mathfrak{g}\) admit the following bases: \((e_3, e_4, e_5, e_6, e_7)\) , \((e_4, e_5, e_6, e_7)\) , \((e_5, e_6, e_7)\) , \((e_6, e_7)\) , \((e_7)\) . Show that the centralizer \(\mathfrak{h}\) of \(\mathcal{C}^5\mathfrak{g}\) admits the basis \((e_2, e_3, e_4, e_5, e_6, e_7)\) .
\item
  Let \((e_i')_{1 \leqslant i \leqslant 7}\) be another basis of \(\mathfrak{g}\) such that \(\mathcal{C}^6\mathfrak{g}\) admits the basis \((e_7')\) , \(\mathcal{C}^5\mathfrak{g}\) admits the basis \((e_6', e_7')\) , \ldots, \(\mathfrak{h}\) admits the basis \((e_2', e_3', e_4', e_5', e_6', e_7')\) . Show that
\end{enumerate}

\[
\left[ e _ {i} ^ {\prime}, e _ {j} ^ {\prime} \right] = \alpha_ {i j} ^ {\prime} e _ {i + j} ^ {\prime} + \beta e _ {i + j + 1} ^ {\prime} + \gamma e _ {i + j + 2} ^ {\prime} + \dots + \lambda e _ {7} ^ {\prime},
\]

where

\[
\alpha_ {1 4} ^ {\prime} \alpha_ {2 5} ^ {\prime} \alpha_ {1 6} ^ {\prime - 1} \alpha_ {2 4} ^ {\prime - 1} = \alpha_ {1 4} \alpha_ {2 5} \alpha_ {1 6} ^ {- 1} \alpha_ {2 4} ^ {- 1}.
\]

\begin{enumerate}
\def\labelenumi{(\alph{enumi})}
\setcounter{enumi}{3}
\tightlist
\item
  Deduce that there exists, if K is infinite, an infinity of non-isomorphic nilpotent Lie algebras of dimension 7 over K and that there exist complex nilpotent Lie algebras which cannot be derived from real nilpotent Lie algebras by extending the base field from R to C.
\end{enumerate}

\begin{enumerate}
\def\labelenumi{\arabic{enumi}.}
\setcounter{enumi}{18}
\tightlist
\item
  \begin{enumerate}
  \def\labelenumii{(\alph{enumii})}
  \tightlist
  \item
    Let \(\mathfrak{g}\) be a Lie algebra. Let \(\mathfrak{S}\) be the Lie algebra of derivations of \(\mathfrak{g}\) . Characteristic ideals \(\mathfrak{g}^{[k]}\) of \(\mathfrak{g}\) are defined inductively as follows: \(\mathfrak{g}^{[0]} = \mathfrak{g}\) and \(\mathfrak{g}^{[k+1]}\) is the subspace of \(\mathfrak{g}\) generated by the Dx ( \(\mathrm{D} \in \mathfrak{S}, x \in \mathfrak{g}^{[k]}\) ). The following conditions are equivalent: (1) every derivation of \(\mathfrak{g}\) is nilpotent: (2) \(\mathfrak{g}^{[k]} = \{0\}\) for \(k\) sufficiently large; (3) the holomorph of \(\mathfrak{g}\) is nilpotent. If they are fulfilled, \(\mathfrak{g}\) is said to be characteristically nilpotent. Such an algebra is nilpotent.
  \end{enumerate}
\end{enumerate}

\begin{enumerate}
\def\labelenumi{(\alph{enumi})}
\setcounter{enumi}{1}
\tightlist
\item
  Verify that an 8-dimensional Lie algebra is defined by the following multiplication table:
\end{enumerate}

\[
\begin{array}{r} \left[ x _ {1}, x _ {2} \right] = x _ {3}, \quad \left[ x _ {1}, x _ {3} \right] = x _ {4}, \quad \left[ x _ {1}, x _ {4} \right] = x _ {5}, \quad \left[ x _ {1}, x _ {5} \right] = x _ {6} \\ \left[ x _ {1}, x _ {6} \right] = x _ {8}, \quad \left[ x _ {1}, x _ {7} \right] = x _ {6}, \quad \left[ x _ {2}, x _ {3} \right] = x _ {5}, \quad \left[ x _ {2}, x _ {4} \right] = x _ {6} \\ \left[ x _ {2}, x _ {5} \right] = x _ {7}, \quad \left[ x _ {2}, x _ {6} \right] = 2 x _ {8}, \quad \left[ x _ {3}, x _ {4} \right] = - x _ {7} + x _ {8}, \quad \left[ x _ {3}, x _ {5} \right] = - x _ {8} \end{array}
\]

and \([x_i, x_j] = 0\) for \(i + j > 8\) . Show that

\[
\begin{array}{c c} \mathcal {C} ^ {2} \mathfrak {g} = \sum_ {i = 3} ^ {8} \mathrm{K} x _ {i}, & \mathcal {C} ^ {3} \mathfrak {g} = \sum_ {i = 4} ^ {8} \mathrm{K} x _ {i}, \\ & \mathcal {C} ^ {6} \mathfrak {g} = \mathrm{K} x _ {8}, \quad \mathcal {C} ^ {7} \mathfrak {g} = \{0 \}, \end{array} \quad \begin{array}{c c} \mathcal {C} ^ {4} \mathfrak {g} = \sum_ {i = 5} ^ {8} \mathrm{K} x _ {i}, & \mathcal {C} ^ {5} \mathfrak {g} = \sum_ {i = 6} ^ {8} \mathrm{K} x _ {i}, \\ & [ \mathcal {C} ^ {2} \mathfrak {g}, \mathcal {C} ^ {2} \mathfrak {g} ] = \mathrm{K} x _ {7} + \mathrm{K} x _ {8}, \end{array}
\]

and that the transporter of \(\mathcal{C}^2\mathfrak{g}\) into \(\mathcal{C}^4\mathfrak{g}\) is \(\sum_{i=2}^{\infty}\mathrm{Kx}_i\) . Deduce that every derivation D of \(\mathfrak{g}\) is defined by formulae \(\mathrm{D}x_i = \sum_{j=1}^{8}u_{ij}x_j\) . Then show that \(u_{ii} = 0\) for all \(i\) if the characteristic of K differs from 2, so that \(\mathfrak{g}\) is characteristically nilpotent.

\begin{enumerate}
\def\labelenumi{(\alph{enumi})}
\setcounter{enumi}{2}
\tightlist
\item
  If the characteristic of K is \(\neq2\) , show that the Lie algebra g of (b) is not the derived Lie algebra of any Lie algebra. (If g = Dh, show first that h is nilpotent. Observe that \(\dim g/Dg = 2\) , which with Exercise 7 (c) leads to a contradiction.)
\end{enumerate}

¶ 20. Let g be a Lie algebra and U its enveloping algebra.

\begin{enumerate}
\def\labelenumi{(\alph{enumi})}
\item
  Let \(\mathfrak{g}'\) be an ideal of \(\mathfrak{g}\) , \(\mathrm{U}' \subset \mathrm{U}\) its enveloping algebra, \(x \in \mathfrak{g}\) and \(a_1, \ldots, a_n\) in U. Suppose that \(a_i = x\) for the indices \(j_1, \ldots, j_n\) and \(a_i \in \mathrm{U}'\) for the other indices (let \(k_1, \ldots, k_q\) be these indices, with \(k_1 < k_2 < \cdots < k_q\) ). Then \(a_1 a_2 \ldots a_n - x^p a_{k_1}a_{k_2} \ldots a_{k_q}\) is a sum of terms of the form \(x^p b\) , where \(b \in \mathrm{U}'\) and \(p' < p\) . (Argue by induction on \(p\) .)
\item
  Suppose henceforth that g is nilpotent and K is of characteristic 0. Let g' be an ideal of codimension 1 in g, U' its enveloping algebra, x an element of g not belonging to g' and U₀ (resp. U₀') the subalgebra of U (resp. U') annihilated by a set b of derivations of g (which extend to derivations of U) mapping g into g'. Suppose that U₀ is contained in the centre of U and U₀ ≠ U'. Show that there exist a₁ ∈ U₀', a₂ ∈ U' such that a₁ ≠ 0 and a = xa₁ + a₂ ∈ U₀. (Let xᵐbₘ + xᵐ⁻¹bₘ₋₁ + ⋯ + b₀ ∈ U₀ with bₘ, \ldots, b₀ in U', m \textgreater{} 0, bₘ ≠ 0. Show, using ⟨a⟩, that, for every derivation D ∈ b of g, Dbₘ = 0, m(Dx)bₘ + Dbₘ₋₁ = 0 and hence D(mxbₘ + bₘ₋₁) = 0.) Show that U₀ is contained in the algebra K{[}a, a₁⁻¹, U₀'{]} generated by a, a₁⁻¹, U₀' in the field of fractions of U₀ (which can be formed by Corollary 7 to Theorem 1 of § 2). Deduce that the field of fractions of U₀ is the field generated by a and U₀'. Show that a is transcendental over K(U₀').
\item
  Let \(\{0\} = \mathfrak{g}_0 \subset \mathfrak{g}_1 \subset \cdots \subset \mathfrak{g}_n = \mathfrak{g}\) be a sequence of ideals of \(\mathfrak{g}\) of dimensions \(0, 1, \ldots, n\) . Let \(x_j\) be an element of \(\mathfrak{g}_j\) not belonging to \(\mathfrak{g}_{j-1}\) . Let \(j_1 < j_2 < \cdots < j_n\) be the indices \(j\) such that there exists in \(\mathrm{U}^j\) (enveloping algebra of \(\mathfrak{g}_j\) ) an element of the centre \(\mathrm{U}_0\) of \(\mathrm{U}\) not belonging to \(\mathrm{U}^{j-1}\) . By (b) there exists \(a_{j_1} \in \mathrm{U}_0 \cap \mathrm{U}^{j_1}\) and \(a_{j_2} \in \mathrm{U}^{j_2-1}\) such that \(a_{j_1} \neq 0\) and \(a_j = x_j a_{j_1} + a_{j_2} \in \mathrm{U}_0 \cap \mathrm{U}^{j}\) . Show by induction on \(n\) that
\end{enumerate}

\[
\mathrm{U} _ {0} \subset \mathrm{K} [ a _ {j _ {1}}, \dots , a _ {j _ {q}}, a _ {j _ {1}} ^ {- 1}, \dots , a _ {j _ {q}} ^ {- 1} ]
\]

and that the field of fractions of \(U_{0}\) is generated by the algebraically independent elements \(a_{j_{1}},\ldots,a_{j_{q}}\) . In particular, the field of fractions of the centre of U is a pure transcendental extension of K.

\begin{enumerate}
\def\labelenumi{\arabic{enumi}.}
\setcounter{enumi}{20}
\tightlist
\item
  Let g be a Lie algebra and σ an automorphism of g.
\end{enumerate}

\begin{enumerate}
\def\labelenumi{(\alph{enumi})}
\tightlist
\item
  Suppose first that \(\mathbf{K}\) is algebraically closed; for all \(\lambda \in \mathbf{K}\) let \(\mathrm{L}_{\lambda}\) be the set of \(x \in \mathfrak{g}\) which are annihilated by a power of \(\sigma - \lambda \mathrm{I}\) ; \(\mathfrak{g}\) is the direct sum of the \(\mathrm{L}_{\lambda}\) . Show that \([\mathrm{L}_{\lambda}, \mathrm{L}_{\mu}] \subset \mathrm{L}_{\lambda \mu}\) . (Note that
\end{enumerate}

\[
(\sigma - \lambda \mu \mathrm{I}) ([ x, y ]) = [ (\sigma - \lambda \mathrm{I}) x, \sigma y ] + [ \lambda x, (\sigma - \mu \mathrm{I}) y ].)
\]

\begin{enumerate}
\def\labelenumi{(\alph{enumi})}
\setcounter{enumi}{1}
\item
  With an arbitrary field K, suppose that none of the eigenvalues of \(\sigma\) (in an algebraically closed extension of K) is a root of unity. Show that g is nilpotent. (Reduce it to the case where K is algebraically closed. If \(\lambda_{1},\ldots,\lambda_{m}\) are the distinct eigenvalues of \(\sigma\) , the \(\lambda_{i}\lambda_{j},\lambda_{i}\lambda_{j}^{2},\lambda_{i}\lambda_{j}^{3},\ldots\) are not all eigenvalues and hence \(ad_{g}x\) is nilpotent for \(x\in L_{\lambda_{i}}\) . Conclude using Exercise 11 applied to the set of \(ad_{g}x\) , where x runs through the union of the \(L_{\lambda_{i}}\) ).
\item
  With an arbitrary field K, suppose that \(\sigma^{q}=I\) , where q is a prime number and that no eigenvalue of \(\sigma\) is equal to 1. Show that g is nilpotent. (Same method as in (b), observing that, if \(g\neq\{0\}\) , q is not equal to the characteristic of K and that for every ordered pair of eigenvalues \(\lambda_{i},\lambda_{j}\) of \(\sigma\) there exists an integer k such that \(\lambda_{i}\lambda_{j}^{k}=1\) .)
\end{enumerate}

\begin{enumerate}
\def\labelenumi{\arabic{enumi}.}
\setcounter{enumi}{21}
\tightlist
\item
  \begin{enumerate}
  \def\labelenumii{(\alph{enumii})}
  \tightlist
  \item
    Let \(u, v\) be two endomorphisms of a vector space \(\mathbf{E}\) of finite dimension over an algebraically closed field \(\mathbf{K}\) . For every eigenvalue \(\lambda\) of \(u\) let \(\mathrm{E}_{\lambda}\) be the subspace of \(\mathbf{E}\) consisting of the vectors annihilated by a power of \(u - \lambda \mathrm{I}\) . Show that, if \((\operatorname{ad} u)^n v = 0\) , the subspaces \(\mathrm{E}_{\lambda}\) are stable under \(v\) .
  \end{enumerate}
\end{enumerate}

\begin{enumerate}
\def\labelenumi{(\alph{enumi})}
\setcounter{enumi}{1}
\item
  Deduce from (a) that, if g is a nilpotent Lie algebra of endomorphisms of E, E is the direct sum of subspaces \(F_{j}\) ( \(1 \leqslant j \leqslant m\) ) which are stable under g and such that on each \(F_{j}\) the restriction of every element \(u \in g\) can be written \(\lambda_{j}(u)\mathbf{I} + u_{j}\) , where \(\lambda_{j}(u) \in \mathbf{K}\) and \(u_{j}\) is nilpotent.
\item
  If \(\mathbf{K}\) is of characteristic 2, \(\mathrm{E} = \mathrm{K}^2\) and \(\mathfrak{g}\) is nilpotent algebra \(\mathfrak{sl}(2,\mathbf{K})\) , show that \(m = 1\) and that possibly \(\lambda (u + v)\neq \lambda (u) + \lambda (v)\) for two elements \(u,v\) of \(\mathfrak{g}\) . (For a sequel to this exercise, cf.~§ 5, Exercise 12.)
\end{enumerate}

\begin{enumerate}
\def\labelenumi{\arabic{enumi}.}
\setcounter{enumi}{22}
\tightlist
\item
  Let \(\mathfrak{g}\) be a Lie \(p\) -algebra over a perfect field \(\mathbf{K}\) of characteristic \(p > 0\) . \(\mathfrak{g}\) is called \(p\) -unipotent if for all \(x \in \mathfrak{g}\) there exists \(m\) such that \(x^{p^m} = 0\) .
\end{enumerate}

\begin{enumerate}
\def\labelenumi{(\alph{enumi})}
\item
  Show that every \(p\) -unipotent Lie \(p\) -algebra is nilpotent.
\item
  Suppose that \(\mathfrak{g}\) is nilpotent and let \(\mathfrak{h}\) denote the \(p\) -core of the centre of \(\mathfrak{g}\) (§ 1, Exercise 23 (a)). Show that \(\mathfrak{g} / \mathfrak{h}\) is \(p\) -unipotent.
\item
  Let \(\mathfrak{g}\) be the nilpotent Lie \(p\) -algebra with a basis of three elements \(e_1, e_2, e_3\) such that \([e_1, e_2] = [e_1, e_3] = 0\) , \([e_2, e_3] = e_1\) , \(e_1^p = e_1\) , \(e_2^p = e_3^p = 0\) . Then \(\mathfrak{h} = \mathrm{Ke}_1\) , but \(\mathfrak{g}\) is not the direct sum of \(\mathfrak{h}\) and a \(p\) -unipotent \(p\) -subalgebra.
\item
  If g is p-unipotent, show that in the restricted enveloping algebra of g the two-sided ideal generated by g is nilpotent. (Argue by induction on the dimension of g.)
\end{enumerate}

\begin{enumerate}
\def\labelenumi{\arabic{enumi}.}
\setcounter{enumi}{23}
\item
  Suppose that the field K is of characteristic 2. Show that in the nilpotent Lie algebra \(\mathfrak{g}_4\) of Exercise 9 there exists no 2-mapping.
\item
  Let \(\mathfrak{g}\) be a Lie \(p\) -algebra. Show that the largest nilpotent ideal of \(\mathfrak{g}\) is a \(p\) -ideal (cf.~§ 1, Exercise 22).
\item
  Let \(\mathbf{G}\) be a group and let \((\mathbf{H}_n)_{n\geqslant 1}\) be a decreasing sequence of subgroups of \(\mathbf{G}\) ; suppose that \(\mathbf{H}_1 = \mathbf{G}\) and that if we write \((x,y) = xyx^{-1}y^{-1}\) , the relations \(x\in \mathrm{H}_i,y\in \mathrm{H}_j\) imply \((x,y)\in \mathrm{H}_{i + j}\) .
\end{enumerate}

\begin{enumerate}
\def\labelenumi{(\alph{enumi})}
\item
  Let \(G_{i} = H_{i}/H_{i+1}\) ; show that \(G_{i}\) is commutative and that the mapping \(x, y \mapsto (x, y)\) defines on passing to the quotient a Z-bilinear mapping of \(G_{i} \times G_{j}\) into \(G_{i+j}\) .
\item
  Let \(\mathrm{gr}(\mathbf{G}) = \sum_{i=1}^{\infty} \mathbf{G}_i\) and extend by linearity the mappings \(\mathbf{G}_i \times \mathbf{G}_j \to \mathbf{G}_{i+j}\) defined in (a) to a Z-bilinear mapping of \(\mathrm{gr}(\mathbf{G}) \times \mathrm{gr}(\mathbf{G})\) into \(\mathrm{gr}(\mathbf{G})\) . Show that \(\mathrm{gr}(\mathbf{G})\) is a Lie Z-algebra with this mapping (to verify the Jacobi identity, use the following formula:
\end{enumerate}

\[
((x, y), z ^ {y}) \cdot ((y, z), x ^ {z}) \cdot ((z, x), y ^ {x}) = e \quad x, y, z \text {   in   } G,
\]

where \(x^{y}\) denotes \(yxy^{-1}\) and \(e\) denotes the identity element of G).

\begin{enumerate}
\def\labelenumi{(\alph{enumi})}
\setcounter{enumi}{2}
\tightlist
\item
  Suppose that there exists \(n\) such that \(\mathrm{H}_n = \{e\}\) . Show that \(\operatorname{gr}(G)\) is a nilpotent Lie \(\mathbf{Z}\) -algebra.
\end{enumerate}

\begin{enumerate}
\def\labelenumi{\arabic{enumi}.}
\setcounter{enumi}{26}
\tightlist
\item
   Let A be an associative algebra with unit element 1 and let \(\mathrm{A}_0 = \mathrm{A} \supset \mathrm{A}_1 \supset \dots \supset \mathrm{A}_n \supset \dots\) be a decreasing sequence of two-sided ideals of A such that \(\mathrm{A}_i \cdot \mathrm{A}_j \subset \mathrm{A}_{i+j}\) . Let G be a group with identity element \(e\) and let \(f: \mathrm{G} \to \mathrm{A}\) be a mapping such that \(f(e) = 1\) , \(f(xy) = f(x) \cdot f(y)\) and \(1 - f(x) \in \mathrm{A}_1\) for all \(x \in \mathrm{G}\) . Let \(\mathrm{H}_n\) denote the set of \(x \in \mathrm{G}\) such that \(1 - f(x) \in \mathrm{A}_n\) . Show that the \(\mathrm{H}_n\) satisfy the conditions of Exercise 26. Show that the mapping \(x \mapsto f(x) - 1\) defines on taking quotients an injective homomorphism of the Lie algebra \(\mathrm{gr}(\mathrm{G})\) into the Lie algebra associated with the graded ring \(\mathrm{gr}(\mathrm{A}) = \sum \mathrm{A}_n / \mathrm{A}_{n+1}\) (cf.~Commutative Algebra, Chapter III).
\end{enumerate}

